% http://web.archive.org/web/20071226183533/http://www.interklasa.pl/euklides/ksiegi.php?str=k02

\chapter{Elementy Euklidesa}
\section{Księga pierwsza, o geometrii płaszczyzny}
\subsection{Księga pierwsza -- definicje}
% Definicja 1.
\begin{euclid_definition}
    Punkt to jest to, co nie składa się z części.
\end{euclid_definition}

% Definicja 2.
\begin{euclid_definition}
    Linia jest długością bez szerokości.
\end{euclid_definition}

% Definicja 3.
\begin{euclid_definition}
    Końcami linii są punkty.
\end{euclid_definition}

% Definicja 4.
\begin{euclid_definition}
    Linia jest prosta, jeżeli położona jest między swoimi punktami w równym i jednostajnym kierunku.
\end{euclid_definition}

% Definicja 5.
\begin{euclid_definition}
    Powierzchnia jest to, co ma tylko długość i szerokość.
\end{euclid_definition}

% Definicja 6.
\begin{euclid_definition}
    Krawędzie powierzchni są liniami.
\end{euclid_definition}

% Definicja 7.
\begin{euclid_definition}
    Płaska powierzchnia albo płaszczyzna jest ta, na której biorąc gdziekolwiek dwa punkty linia prosta między tymi punktami cała leży na tej powierzchni.
\end{euclid_definition}

% Definicja 8.
\begin{euclid_definition}
    Kąt płaski to nachylenie dwóch linii na płaszczyźnie w miejscu, w którym jedna spotyka drugą i nie leżą w linii prostej.
\end{euclid_definition}

% Definicja 9.
\begin{euclid_definition}
    Kiedy linie są proste i tworzą kąt, wtedy kąt zwany jest prostoliniowym.
\end{euclid_definition}

% Definicja 10.
\begin{euclid_definition}
    Kiedy linia prosta padająca na drugą linie prostą, tworzy z nią kąty przyległe równe między sobą, to każdy z kątów równych nazywamy prostym, a padająca linia prostą nazywa się prostopadłą do tej linii, na którą pada.
\end{euclid_definition}

% Definicja 11.
\begin{euclid_definition}
    Kąt rozwarty jest większy od kąta prostego.
\end{euclid_definition}

% Definicja 12.
\begin{euclid_definition}
    Kąt ostry jest mniejszy od kąta prostego.
\end{euclid_definition}

% Definicja 13.
\begin{euclid_definition}
    Kresem albo granicą jest to, na czym się dana rzecz kończy.
\end{euclid_definition}

% Definicja 14.
\begin{euclid_definition}
    Figurą nazywamy to co jest ograniczone granicą lub granicami. 
\end{euclid_definition}

% Definicja 15.
\begin{euclid_definition}
    Koło jest figurą płaską zawarta linią zwaną okręgiem, do której wszystkie linie proste poprowadzone z jednego punktu wewnątrz figury położonego, są między sobą równe.
\end{euclid_definition}

% Definicja 16.
\begin{euclid_definition}
    I ten punkt nazywa się centrum lub środkiem koła. 
\end{euclid_definition}

% Definicja 17.
\begin{euclid_definition}
    Średnicą koła jest każda linia narysowana przez środek koła, przedłużona w dwóch kierunkach do jego obwodu, przepoławiająca go.
\end{euclid_definition}

% Definicja 18.
\begin{euclid_definition}
    Półokręgiem jest figura zawarta między średnicą i częscia okręgu odciętą tą średnicą. Środek półokregu jest też środkiem okręgu.
\end{euclid_definition}

% Definicja 19.
\begin{euclid_definition}
    Figury prostokreślne to figury ograniczone prostymi. Trójkąt to figura prostokreślna ograniczona trzema prostymi. Czworobok lub czworokąt to figura prostokreślna, która jest ograniczona czterema prostymi. Wielobok lub wielokąt to figura prostokreślna ograniczona więcej niż czterema prostymi.
\end{euclid_definition}

% Definicja 20.
\begin{euclid_definition}
    Trójkąt równoboczny to trójkąt, który ma trzy boki równe. Trójkąt równoramienny to trójkąt, który ma tylko dwa boki równe. Trójkąt różnoboczny to trójkąt, który ma trzy boki różne.
\end{euclid_definition}

% Definicja 21.
\begin{euclid_definition}
    Ponadto: trójkąt prostokątny to trójkąt, który na kąt prosty. Trójkąt rozwartokątny to trójkąt, który ma kąt rozwarty. Trójkąt ostrokątny to trójkąt, który ma trzy kąty ostre.
\end{euclid_definition}

% Definicja 22.
\begin{euclid_definition}
    Kwadrat jest to czworobok mający równe boki i równe kąty. Prostokąt jest to czworobok mający kąty proste, ale boki nierówne. Romb (kwadrat ukośny) jest to czworobok mający równe boki, ale nie mający kątów prostych. Równoległobok jest to czworobok mający boki przeciwległe równe, ale nie mający katów prostych. Wszystkie czworoboki inne niż wyżej wymienione nazywamy czworokątami.
\end{euclid_definition}

% Definicja 23.
\begin{euclid_definition}
    Linie równoległe, czyli mówiąc krócej równoległe są to proste, które leżą na tej samej płaszczyźnie i przedłużone z obu stron w nieskończoność, z żadnej strony nie przetną się.
\end{euclid_definition}
\subsection{Księga pierwsza -- twierdzenia}
\setcounter{counter_euclid}{0}

% Twierdzenie 1.
\begin{euclid_theorem}
    Na danej linii prostej skonstruuj trójkąt równoboczny o żadanych bokach.
\end{euclid_theorem}

% Twierdzenie 2.
\begin{euclid_theorem}
    Skonstruuj odcinek równy danemu odcinkowi którego koniec jest zadanym punktem.
\end{euclid_theorem}

% Twierdzenie 3.
\begin{euclid_theorem}
    Mając dane dwie linie proste nierówne, od większej odciąć linię równą mniejszej.
\end{euclid_theorem}

% Twierdzenie 4.
\begin{euclid_theorem}
    Jeśli dwa trójkąty mają dwa boki odpowiednio równe dwóm innym, i jeżeli kąty zawarte między bokami równoległymi są równe, wtedy ich podstawy również są sobie równe i pozostałe kąty równe są odpowiednim kątom.
\end{euclid_theorem}

% Twierdzenie 5.
\begin{euclid_theorem}
    W trójkątach równoramiennych kąty przy podstawie są sobie równe oraz kąty powstałe przez przedłużenie boków równych są sobie równe.
\end{euclid_theorem}

% Twierdzenie 6.
\begin{euclid_theorem}
    Jeżeli w trójkącie dwa kąty są sobie równe, to boki leżące naprzeciw równych kątów są sobie równe.
\end{euclid_theorem}

% Twierdzenie 7.
\begin{euclid_theorem}
    Na tej samej podstawie i z tej samej strony nie mogą być wykreślone dwa trójkąty takie, żeby boki w tych trójkątach przy obydwu końcach wspólnej podstawy były między sobą równe.
\end{euclid_theorem}

% Twierdzenie 8.
\begin{euclid_theorem}
    Jeżeli dwa boki jednego trójkąta są równe dwóm bokom drugiego trójkąta, to kąty zawarte między równymi bokami są sobie równe.
\end{euclid_theorem}

% Twierdzenie 9.
\begin{euclid_theorem}
    Dany kąt prostokreślny podziel na dwie równe części.
\end{euclid_theorem}

% Twierdzenie 10.
\begin{euclid_theorem}
    Daną linię prostą oznaczoną podzielić na dwie równe części.
\end{euclid_theorem}

% Twierdzenie 11.
\begin{euclid_theorem}
    Z punktu danego na danej linii prostej wyprowadzić linie prostopadłą do danej linii prostej.
\end{euclid_theorem}

% Twierdzenie 12.
\begin{euclid_theorem}
    Z punktu danego leżącego poza linią prostą nieograniczoną, wyprowadzić prostą linię prostopadłą do niej.
\end{euclid_theorem}

% Twierdzenie 13.
\begin{euclid_theorem}
    Jeżeli linia prosta przecinająca drugą prostą tworzy z nią dwa kąty, to są one proste, albo równe dwóm kątom prostym.
\end{euclid_theorem}

% Twierdzenie 14.
\begin{euclid_theorem}
    Jeżeli przy linii prostej i przy punkcie na niej leżącym dwie linie proste nie po jednej stronie położone czynią kąty przyległe równe dwóm kątom prostym, to te linie proste będą w tym samym kierunku.
\end{euclid_theorem}

% Twierdzenie 15.
\begin{euclid_theorem}
    Jeżeli dwie linie proste przecinają się, to utworzone przez nie kąty przeciwległe są sobie równe.
\end{euclid_theorem}

% Twierdzenie 16.
\begin{euclid_theorem}
    W dowolnym trójkącie kąt zewnętrzny powstały przez przedłużenie jednego boku jest większy od każdego z dwóch kątów wewnętrznych przeciwległych jemu.
\end{euclid_theorem}

% Twierdzenie 17.
\begin{euclid_theorem}
    W każdym trójkącie suma dwóch dowolnych kątów jest mniejsza od dwóch kątów prostych.
\end{euclid_theorem}

% Twierdzenie 18.
\begin{euclid_theorem}
    W każdym trójkącie bok większy przeciwległy jest kątowi większemu.
\end{euclid_theorem}

% Twierdzenie 19.
\begin{euclid_theorem}
    W każdym trójkącie kąt większy przeciwległy jest bokowi większemu.
\end{euclid_theorem}

% Twierdzenie 20.
\begin{euclid_theorem}
    W każdym trójkącie suma dwóch dowolnych boków jest większa od boku trzeciego.
\end{euclid_theorem}

% Twierdzenie 21.
\begin{euclid_theorem}
    Jeżeli z końców jednego boku trójkąta poprowadzone będą dwie linie proste wewnątrz trójkąta, aż do zejścia się z sobą, to te dwie linie proste będą mniejsze od dwóch pozostałych boków trójkąta, lecz zawierać jednak 
\end{euclid_theorem}
będą kąt większy od kąta zawartego między pozostałymi bokami trójkąta.
% Twierdzenie 22.
\begin{euclid_theorem}
    Aby z trzech danych linii prostych wykreślić trójkąt, potrzeba aby z tych trzech danych linii prostych suma dwóch którychkolwiek była większa od trzeciej.
\end{euclid_theorem}

% Twierdzenie 23.
\begin{euclid_theorem}
    Na danej linii prostej i punkcie na niej danym wykreślić kąt prostokreślny równy kątowi prostokreślnemu danemu.
\end{euclid_theorem}

% Twierdzenie 24.
\begin{euclid_theorem}
    Jeżeli dwa boki jednego trójkąta, są równe dwóm bokom trójkąta drugiego, z kątów zaś między bokami równymi jeden większy jest od drugiego; to będzie też podstawa jednego trójkąta większa od podstawy drugiego trójkąta.
\end{euclid_theorem}

% Twierdzenie 25.
\begin{euclid_theorem}
    Jeżeli dwa boki jednego trójkąta, są równe dwóm bokom trójkąta drugiego, lecz podstawa jednego trójkąta większa jest od podstawy drugiego trójkąta, to i kąty między bokami równymi zawarte będą jeden większy od drugiego.
\end{euclid_theorem}

% Twierdzenie 26.
\begin{euclid_theorem}
    Jeżeli dwa kąty jednego trójkąta są równe dwóm kątom drugiego trójkąta, i bok jeden przyległy obydwu kątom, albo jednemu w pierwszym trójkącie równa się bokowi jednemu przyległemu obydwu katom, albo jednemu w drugim trójkącie; będą i dwa boki pozostałe równe dwóm bokom pozostałym i kąt trzeci w jednym trójkącie będzie równy katowi trzeciemu w drugim trójkącie.
\end{euclid_theorem}

% Twierdzenie 27.
\begin{euclid_theorem}
    Jeżeli na dwie linie proste, pada linia prosta czyniąca kąty naprzemian równe między sobą, to te dwie linie proste będą równoległe.
\end{euclid_theorem}

% Twierdzenie 28.
\begin{euclid_theorem}
    Jeśli linia prosta opada na dwie linie proste, tworząc kąt zewnętrzny równy wewnętrznemu i przeciwny do kąta na tym samym boku lub suma kątów wewnętrznych na tym samym boku jest równa dwóm kątom prostym, wtedy linie proste są równoległe do siebie.
\end{euclid_theorem}

% Twierdzenie 29.
\begin{euclid_theorem}
    Linia prosta opada na równoległą linie prostą tworząc alternatywne kąty równe sobie, kąt zewnętrzny równy wewnętrznemu i przeciwległy i suma kątów wewnętrznych na tym samym boku jest równa dwóm kątom prostym.
\end{euclid_theorem}

% Twierdzenie 30.
\begin{euclid_theorem}
    Linie proste, które są równoległe do linii prostej są również równoległe do siebie.
\end{euclid_theorem}

% Twierdzenie 31.
\begin{euclid_theorem}
    Poprowadzić przez dany punkt linię prostą równoległą względem danej lini prostej.
\end{euclid_theorem}

% Twierdzenie 32.
\begin{euclid_theorem}
    W jakimkolwiek trójkącie, jeśli jeden z boków jest znany wtedy kąt zewnętrzny jest równy sumie dwóch kątów wewnętrznych i przeciwnych i suma trzech wewnętrznych kątów trójkąta jest równa dwóm kątom prostym.
\end{euclid_theorem}

% Twierdzenie 33.
\begin{euclid_theorem}
    Linie proste, które łączą końce równych i równoległych linii prostych w tym samym kierunku są sobie równe i równoległe.
\end{euclid_theorem}

% Twierdzenie 34.
\begin{euclid_theorem}
    W równoległobokach boki i kąty przeciwne są między sobą równe, a przekątna dzieli je na dwie równe części.
\end{euclid_theorem}

% Twierdzenie 35.
\begin{euclid_theorem}
    Równoległoboki, które są na takiej samej podstawie i są porównywalne są sobie równe.
\end{euclid_theorem}

% Twierdzenie 36.
\begin{euclid_theorem}
    Równoległoboki, które mają równe podstawy i są porównywalne są sobie równe.
\end{euclid_theorem}

% Twierdzenie 37.
\begin{euclid_theorem}
    Trójkąty, które mają takie same podstawy i są porównywalne są sobie równe.
\end{euclid_theorem}

% Twierdzenie 38.
\begin{euclid_theorem}
    Trójkąty, których podstawy są równe i są one porównywalne są sobie równe.
\end{euclid_theorem}

% Twierdzenie 39.
\begin{euclid_theorem}
    Równe trójkąty, które są na takich samych podstawach i mające te same boki również są porównywalne.
\end{euclid_theorem}

% Twierdzenie 40.
\begin{euclid_theorem}
    Równe trójkąty, które mają takie same podstawy i mają te same boki również są porównywalne.
\end{euclid_theorem}

% Twierdzenie 41.
\begin{euclid_theorem}
    Jeśli równoległobok i trójkąt mają tą samą podstawę i są tymi samymi liniami zakończone, to trójkąt jest połową równoległoboku.
\end{euclid_theorem}

% Twierdzenie 42.
\begin{euclid_theorem}
    Skonstruować równoległobok równy danemu trójkątowi o podanym prostoliniowym kącie.
\end{euclid_theorem}

% Twierdzenie 43.
\begin{euclid_theorem}
    W każdym równoległoboku, dopełnienia równoległoboków koło przekątnych położonych są między sobą równe.
\end{euclid_theorem}

% Twierdzenie 44.
\begin{euclid_theorem}
    Na danej linii prostej wykreślić równy danemu równoległobok, którego jeden kąt będzie równy danemu.
\end{euclid_theorem}

% Twierdzenie 45.
\begin{euclid_theorem}
    Wykreślić równy danej figurze prostokreślny równoległobok, którego jeden kąt będzie równy danemu.
\end{euclid_theorem}

% Twierdzenie 46.
\begin{euclid_theorem}
    Na danej linii prostej wykreślić kwadrat.
\end{euclid_theorem}

% Twierdzenie 47.
\begin{euclid_theorem}
    W trójkącie prostokątnym, kwadrat zbudowany na boku przeciwnym kątowi prostemu, równy jest kwadratom zbudowanym na bokach, które kąt prosty zawierają.
\end{euclid_theorem}

% Twierdzenie 48.
\begin{euclid_theorem}
    Jeżeli kwadrat zbudowany na jednym z boków trójkąta, jest równy kwadratom wykreślonym na dwóch pozostałych bokach trójkąta, to kąt zawarty między dwoma pozostałymi bokami będzie prosty.
\end{euclid_theorem}
\subsection{Księga pierwsza -- postulaty}
\setcounter{counter_euclid}{0}

% Postulat 1
\begin{euclid_postulate}
    Można poprowadzić prostą od któregokolwiek punktu do któregokolwiek punktu.
\end{euclid_postulate}

% Postulat 2.
\begin{euclid_postulate}
    Ograniczoną prostą można przedłużyć nieskończenie.
\end{euclid_postulate}

% Postulat 3.
\begin{euclid_postulate}
    Można zakreślić okrąg z któregokolwiek punktu jako środka dowolną odległością.
\end{euclid_postulate}

% Postulat 4.
\begin{euclid_postulate}
    Wszystkie kąty proste są między sobą równe.
\end{euclid_postulate}

% Postulat 5.
\begin{euclid_postulate}
    Jeżeli prosta przecinająca dwie proste tworzy z nimi kąty jednostronnie wewnętrzne o sumie mniejszej niż dwa kąty proste, to te dwie proste przedłużone nieskończenie przecinają się po tej stronie, po której znajdują się kąty o sumie mniejszej od dwóch kątów prostych.
\end{euclid_postulate}
\subsection{Księga pierwsza -- pojęcia podstawowe}
\setcounter{counter_euclid}{0}
% Pojęcie podstawowe 1
\begin{euclid_basic}
Wyrażenia, które są równe się temu samemu wyrażeniowi, są sobie równe.
\end{euclid_basic}

% Pojęcie podstawowe 2
\begin{euclid_basic}
Jeżeli równania dodawane są do równań, wtedy całości są sobie równe.
\end{euclid_basic}

% Pojęcie podstawowe 3
\begin{euclid_basic}
Jeżeli równania odejmowane są do równań, wtedy całości są sobie równe.
\end{euclid_basic}

% Pojęcie podstawowe 4
\begin{euclid_basic}
Wyrażenia, które się pokrywają, są sobie równe.
\end{euclid_basic}

% Pojęcie podstawowe 5
\begin{euclid_basic}
Całość jest większa od części. 
\end{euclid_basic}













\section{Księga druga, o Iksach Dupiksach}
    Księga II - Definicje:

Definicja 1
    Każdy równoległobok prostokątny wyraża i wykreśla się dwiema liniami prostymi które zawierają właściwy kąt.
Definicja 2
    W równoległoboku jeżeli poprowadzimy przekątną i przez punkt gdziekolwiek obrany na tej przekątnej poprowadzimy dwie linie równoległe do boków równoległoboku, równoległobok podzieli się na cztery części, każda z dwóch części której przekątna jest częścią przekątnej całego równoległoboku, wzięta z dwiema jej przyległymi zwać będziemy węgielnicą.


Księga II - Twierdzenia:

Twierdzenie 1
    Jeżeli z dwóch linii prostych podzielimy jedną którąkolwiek na ilekolwiek części (które będziemy nazywać odcinkami), prostokąt zawarty dwiema liniami prostymi, równy będzie prostokątom wykreślonym z linii prostej nieprzecietej i z odcinków drugiej linii prostej.
Twierdzenie 2
    Jeżeli linie prostą podzielimy jakkolwiek, prostokąty zawarte całą linią i jej oddzielnymi odcinkami będą równe kwadratowi z całej linii.
Twierdzenie 3
    Jeżeli linie prostą podzielimy na dwa jakiekolwiek odcinki; prostokąt całą linią i jednym odcinkiem, zawarty, będzie równy prostokątowi odcinkami linii prostej zawartymi wraz z kwadratem wyrażonym na odcinku wziętym z boku drugiego prostokąta pierwszego.
Twierdzenie 4
    Jeżeli linię prostą podzielimy na dwa jakiekolwiek odcinki, kwadrat z całej linii będzie równy kwadratom z obydwu odcinków linii dwa razy wziętemu prostokątowi zawartemu odcinkami linii.
Twierdzenie 5
    Jeżeli linię prostą podzielimy na dwa równe odcinki i na dwa odcinki nierówne; to prostokąt odcinkami nierównymi zawartymi wraz z kwadratem wystawionym na odcinkach między podziałami zawartymi będzie równy kwadratowi wystawionemu na połowie linii.
Twierdzenie 6
    Jeżeli linię prostą na dwa różne odcinki podzieloną przedłużymy podług upodobani; prostokąt zawarty linią prostą wraz z przedłużeniem wziętą i samym przedłużeniem, wraz z kwadratem wystawionym na połowie linii, będzie równy kwadratowi wystawionemu na połowie linii wraz z przedłużeniem wziętym.
Twierdzenie 7
    Jeżeli linię prostą podzielimy na dwa różne odcinki nierówne; kwadraty: pierwszy z całej linii, drugi z jej odcinka, będą równe dwa razy wziętemu prostokątowi całą linią i tym samym odcinkiem zawartym wraz z kwadratem z odcinka drugiego.
Twierdzenie 8
    Jeżeli linię podzielimy na dwa odcinki nierówne; cztery razy wzięty prostokąt całą linią i jej jednym odcinkiem zawarty wraz z kwadratem z odcinka drugiego, będzie równy kwadratowi wystawionemu na linii złożonej z całej linii i z odcinka pierwszego. 
Twierdzenie 9
    Jeżeli linię prostą podzielimy na dwa odcinki równe, i na dwa odcinki nierówne; kwadraty z odcinków nierównych będą dwa razy większe od kwadratów, z których jeden byłby wystawiony na połowie linii, drugi na linii miedzy podziałami zawartej.
Twierdzenie 10
    Jeżeli linię prostą na dwa odcinki równe podzieloną przedłużymy według upodobania; kwadraty: pierwszy z całej linii wraz z przedłużeniem, drugi z samego przedłużenia, będą dwa razy większe od kwadratów, z których pierwszy byłby wystawiony na połowie linii, a drugi na połowie linii wraz z przedłużeniem wziętym.
Twierdzenie 11
    Daną linię prostą podzielić na dwa odcinki tak, aby prostokąt całą linią i jednym jej odcinkiem zawarty, był równy kwadratowi z odcinka drugiego.
Twierdzenie 12
    W trójkątach rozwartokątnych, kwadrat z boku kątowi przeciwnemu rozwartemu, większy jest od kwadratów z ramion kąta rozwartego o dwa razy wzięty prostokąt, zawarty ramionami kąta rozwartego i przedłużeniem tego ramienia zamkniętym między wierzchołkiem kąta rozwartego i punktem w którym prostopadła z końca drugiego ramienia kąta rozwartego spuszczona na pierwsze ramie, spotyka przedłużenie odcinka.
Twierdzenie 13
    W każdym trójkącie, kwadrat z boku przeciwnego kątowi ostremu, mniejszy jest od kwadratów z ramion ten kąt obejmujących, o dwa razy wzięty prostokąt zawarty ramieniem tego kąta ostrego i odcinka, lub przedłużeniem tego ramienia zamkniętym między wierzchołkami kata ostrego i punktem, w którym linia prostopadła z końca drugiego ramienia kąta ostrego spuszczona na pierwsze ramię spotyka to ramię lub przedłużenie danego ramienia.
Twierdzenie 14
    Na danej figurze prostokreślnej równy kwadrat wykreślić. 

\section{Księga trzecia, o Iksach Dupiksach}
    Księga III - Definicje:

Definicja 1.
    Koła równe są to koła których średnice lub promienie są równe.
Definicja 2.
    Mówi się, że linia prosta dotyka koła, gdy będąc styczną z kołem przedłużona z obydwu stron nie przecina się z żadnej strony okręgu koła.
Definicja 3.
    Mówi się, że koła dotykają się wzajemnie, gdy te koła prócz jednego punktu stykania się ich okręgów, same się nie przecinają.
Definicja 4.
    Mówi się, że linie proste równoodległe są od środka koła, gdy prostopadłe ze środka koła na nie spuszczone są równe.
Definicja 5.
    Mówi się, że ta linia prosta bardziej jest odległa od środka koła, na którą prostopadła ze środka koła spuszczona jest większa.
Definicja 6.
    Odcinkiem koła jest figura czyli część koła ograniczona linią prostą i okręgiem koła.
Definicja 7.
    Kąt zaś odcinka jest ten, który się linią prostą i okręgiem koła zawiera.
Definicja 8.
    Jeżeli na okręgu koła wzięty będzie punkt i od niego będą poprowadzone linie proste do końców linii prostej za podstawę odcinkami służącej, kąt między tymi liniami prostymi zawarty jest kątem w odcinku.
Definicja 9.
    Kiedy zaś linie proste kąt zawierające zajmują część okręgu, mówi się, że kąt ten opiera się na okręgu koła.
Definicja 10.
    Jeżeli kąt ma swój wierzchołek we środku koła; figura czyli część koła zawarta między ramionami tegoż koła, to jest między promieniami i łukiem koła nazywa się wycinkiem koła.
Definicja 11.
    Odcinkami podobnymi kół nazywają się te, które zajmują kąty równe, lub w których kąty są równe między sobą.


Księga III - Twierdzenia:

Twierdzenie 1.
    Wyznaczyć środek koła danego.
Twierdzenie 2.
    Jeżeli na okręgu obierzemy dwa gdziekolwiek punkty, linia prosta łącząca te punkty padnie wewnątrz koła.
Twierdzenie 3.
    Jeżeli w kole linia prosta przez środek poprowadzona przecina linie nie przez środek poprowadzoną na dwie równe części, będzie pierwsza prostopadła do drugiej; i jeżeli pierwsza jest prostopadła do drugiej, przecina ja na dwie równe części.
Twierdzenie 4.
    Jeżeli w kole dwie linie proste, nie przez środek koła poprowadzone przecinają się nawzajem, nie przetną się na dwie równe części. 
Twierdzenie 5.
    Jeżeli dwa koła przecinają się nawzajem, wspólnego środka mieć nie mogą.
Twierdzenie 6.
    Jeżeli dwa koła dotykają się wzajemnie, to wspólnego środka mieć nie mogą.
Twierdzenie 7.
    Jeżeli na średnicy koła wzięty będzie punkt którykolwiek oprócz średnicy koła i od tego punktu poprowadzone linie proste do okręgu, ze wszystkich linii największa będzie część średnicy, na której znajduje się środek koła, a najmniejsza pozostała część średnicy, z innych zaś linii prostych każda bliższa przechodząca przez środek koła, większa będzie od odleglejszej, z tego na koniec punktu dwie tylko równe linie proste z obydwu stron najmniejszej linii prostej mogą być do okręgu poprowadzone.
Twierdzenie 8.
    Jeżeli z punktu zewnątrz koła obranego, poprowadzone będą do okręgu linie proste, z których jedna przechodziła by przez środek koła a inne padały gdziekolwiek, z linii prostych padających na część okręgu wklęsłą, największa jest linia poprowadzona przez środek koła, z innych zaś linii każda bliższa przechodzącej przez środek jest większa od odleglejszej. Lecz z linii padających na cześć okręgu wypukłą, najmniejsza jest linia prosta zawarta między punktem zewnętrz koła i średnicą, z innych zaś linii prostych każda bliższa najmniejszej, mniejsza jest odleglejsza; na koniec dwie tylko równe linie proste z tego punktu po obydwu stronach najmniejszej linii prostej mogą być do okręgu poprowadzone.
Twierdzenie 9.
    Jeżeli z punktu danego wewnątrz koła poprowadzimy do okręgu więcej niż dwie linie proste i te proste są miedzy sobą równe, punkt ten będzie środkiem koła.
Twierdzenie 10.
    Gdy koła przecinają się to okręgi ich w dwóch tylko punktach się przecinają.
Twierdzenie 11.
    Jeżeli dwa koła stykają się ze sobą wewnątrz, linia łącząca środki tychże kół przedłużona pada na punkt dotykania się kół.
Twierdzenie 12.
    Jeżeli dwa koła dotykają się ze sobą zewnętrznie, to linia prosta łącząca ich środki przechodzi przez punkt dotykania się.
Twierdzenie 13.
    Okrąg koła nie może dotykać okręgu drugiego koła w więcej niż jednym punkcie, nieważne jest czy dotkniecie jest zewnętrzne bądź wewnętrzne.
Twierdzenie 14.
    W kole linie proste równe, na okręgu jego zakończone, są równoodległe od środka; i linie proste które na okręgu jego zakończone są równoodległe od środka, są też miedzy sobą równe.
Twierdzenie 15.
    Ze wszystkich linii prostych w kole poprowadzonych i na okręgu jego zakończonych, największa jest średnica, z innych zaś każda bliższa środka koła, większa jest od odleglejszej; i z dwóch linii prostych nierównych, większa bliższa jest środka koła od mniejszej.
Twierdzenie 16.
    Prostopadła do średnicy koła z końca jej wyprowadzona, pada cała zewnątrz koła, a między tą prostopadłą i okręgiem żadna inna linia prosta nie pada; albo tak samo: okrąg koła przechodzi miedzy prostopadłą do średnicy i linią prostą, która ze średnicą kąt ostry jakokolwiek wielki zawiera, czyli która zawiera kąt jakokolwiek mały z prostopadłą do średnicy.
Twierdzenie 17.
    Z danego punktu za kołem, lub na jego okręgu wyprowadzić linię prostą, która by danego koła dotykała.
Twierdzenie 18.
    Jeżeli linia prosta dotyka się okręgu koła, a ze środka koła wyprowadzona będzie linia prosta do punktu dotykania się, to ta będzie prostopadła do stycznej.
Twierdzenie 19.
    Jeżeli linia prosta dotyka okręgu koła, z punktu zaś dotknięcia wyprowadzona będzie do tej stycznej prostopadła, to na prostopadłej będzie środek koła.
Twierdzenie 20.
    W kole, kąt mający wierzchołek we środku jest podwojeniem kata mającego swój wierzchołek na okręgu koła, gdyż tę samą podstawę okręgu mają za podstawę, czyli to samo gdy ramionami swymi tej samej części okręgu obejmują.
Twierdzenie 21.
    Kąty w tym samym odcinku koła są między sobą równe.
Twierdzenie 22.
    Kąty przeciwne czworokąta w koło wpisane są równe dwóm kątom prostym.
Twierdzenie 23.
    Na tej samej linii prostej nie można wykreślić dwóch odcinków kół po tej samej stronie podobnych, które by nie przystawały do siebie.
Twierdzenie 24.
    Wykreślone na równych liniach prostych podobne odcinki kół, są między sobą równe.
Twierdzenie 25.
    Mając dany odcinek koła, opisać koła którego jest odcinkiem.
Twierdzenie 26.
    W kołach równych, kąty równe w środkach lub przy okręgach wspierają się na równych łukach.
Twierdzenie 27.
    W kołach równych, kąty we środkach lub przy okręgach, na równych łukach wspierające się, są między sobą równe.
Twierdzenie 28.
    W kołach równych, cięciwy równe obejmują łuki równe, tak, że łuk większy większemu, mniejszy mniejszemu jest równy.
Twierdzenie 29.
    W kołach równych, równe łuki obejmują cięciwy równe.
Twierdzenie 30.
    Dany łuk podzielić na dwie części.
Twierdzenie 31.
    W kole, kąt w półkolu jest prosty; z katów zaś w odcinkach nierównych, kąt w większym odcinku mniejszy jest od prostego; a w mniejszym odcinku większy od prostego.
Twierdzenie 32.
    Jeżeli okręgu koła dotyka linia prosta, z punktu zaś dotknięcia poprowadzona będzie cięciwa, kąty zawarte miedzy cięciwową i styczną, będą równe kątom w odcinkach koła na przemian.
Twierdzenie 33.
    Na danej linii prostej wykreślić odcinek koła który by zawierał kąt równy kątowi danemu.
Twierdzenie 34.
    Z koła danego oddzielić odcinek któryby zawierał kąt równy danemu kątowi.
Twierdzenie 35.
    Jeżeli w kole dwie cięciwy przecinają się nawzajem, prostokąt zawarty odcinkami jednej cięciwy będzie równy prostokątowi zawartemu odcinkami drugiej cięciwy.
Twierdzenie 36.
    Jeżeli z punktu za kołem obranego, poprowadzimy dwie linie proste, których jedna przecinałaby koło, a druga byłaby styczną; to prostokąt zawarty całą linia przecinającą i odcinkiem jej za kołem będzie równy kwadratowi ze stycznej.
Twierdzenie 37.
    Jeżeli z dwóch linii prostych, od jednego punktu zewnątrz koła obranego poprowadzonych, jedna przecina koło, a druga pada na okrąg tego koła: i jeżeli prostokąt z całej linii przecinającej i odcinka jej za kołem będącego jest równy kwadratowi z linii padającej na okrąg koła, to linia będzie padająca na okrąg koła styczną. 

\section{Księga piąta, o Iksach Dupiksach}

Definicja 1.
    Wielkość jest częścią innej wielkości, mniejsza większej, jeśli mniejsza mierzy większą.
Definicja 2.
    Większa wielkość jest wielokrotnością mniejszej wielkości, jeśli mniejsza mierzy większą.
Definicja 3.
    Stosunek to rodzaj zależności pomiędzy dwiema wielkościami tego samego typu.
Definicja 4.
    Mówi się, że wielkości są w stosunku między sobą, jeśli jedna z nich zwielokrotniona, może przewyższyć drugą.
Definicja 5.
    Mówi się, że wielkości są w tym samym stosunku, pierwsza do drugiej i trzecia do czwartej, jeśli wziąwszy równe wielokrotności pierwszej i trzeciej oraz równe wielokrotności drugiej i czwartej, otrzymujemy zależność: gdy wielokrotność pierwszej wielkości jest większa, równa lub mniejsza od wielokrotności trzeciej wielkości, to zależność pomiędzy wielokrotnością drugiej i wielokrotnością czwartej jest taka sama.
Definicja 6.
    Wielkości, które są w tym samym stosunku, nazywamy proporcjonalnymi.
Definicja 7.
    Jeśli, z równych wielokrotności, wielokrotność pierwszej przewyższa wielokrotność drugiej, ale wielokrotność trzeciej nie przewyższa wielokrotności czwartej, to mówi się, że pierwsza do drugiej jest w większym stosunku niż trzecia do czwartej.
Definicja 8.
    Proporcja składa się najmniej z trzech wyrazów.
Definicja 9.
    Jeśli trzy wielkości są proporcjonalne, to mówi się, że pierwsza wielkość do trzeciej jest w stosunku podwójnym do stosunku pierwszej i drugiej.
Definicja 10.
    Jeśli cztery wielkości są ciągło proporcjonalne, to mówi się, że pierwsza wielkość do czwartej jest w potrójnym stosunku do stosunku pierwszej i drugiej, i tak dalej, ilekolwiek byłoby wielkości proporcjonalnych.
Definicja 11.
    Mówi się, że poprzedniki są wielkościami odpowiadającymi poprzednikom, a następniki są wielkościami odpowiadającymi następnikom.
Definicja 12.
    Stosunek przemienny oznacza, że, jeśli cztery wielkości są proporcjonalne, to pierwsza wielkość do trzeciej jest w takim stosunku, jak druga do czwartej.
Definicja 13.
    Odwrócony stosunek (inwersja stosunku) oznacza, że, jeśli cztery wielkości są proporcjonalne, to druga do pierwszej ma się tak, jak czwarta do trzeciej.
Definicja 14.
    Stosunek "dodawanie wyrazów" oznacza, że, jeśli cztery wielkości są proporcjonalne, to pierwsza wraz z drugą ma się do drugiej, jak trzecia wraz z czwartą do czwartej.
Definicja 15.
    Stosunek "odejmowanie wyrazów" oznacza, że, jeśli cztery wielkości są proporcjonalne, to nadmiar pierwszej nad drugą ma się do drugiej tak, jak nadmiar trzeciej nad czwartą do czwartej.
Definicja 16.
    Zamiana stosunku oznacza, że, jeśli cztery wielkości są proporcjonalne, to pierwsza ma się do nadmiaru pierwszej nad drugą tak, jak trzecia do nadmiaru trzeciej nad czwartą.
Definicja 17.
    Stosunek "ex aequali" pojawia się, jeśli mamy dwie równoliczne grupy wielkości, które w obrębie swojej grupie są proporcjonalne parami oraz w grupie pierwszych wielkości- pierwsza do ostatniej ma się tak, jak w drugiej grupie wielkości - wielkość pierwsza do ostatniej. Lub ,innymi słowy, oznacza to zamianę skrajnych wyrazów na środkowe.
Definicja 18.
    Proporcja "perturbata" pojawia się, jeśli mamy dwie trzy-elementowe grupy wielkości oraz pierwsza wielkość do drugiej, spośród wielkości z pierwszej grupy, ma się tak, jak druga wielkość do trzeciej, spośród wielkości drugiej grupy oraz druga do trzeciej z pierwszej grupy ma się tak, jak pierwsza do drugiej z drugiej grupy.


Księga V - Twierdzenia:

Twierdzenie 1.
    Jeżeli każda z dowolnej liczby wielkości jest taką samą wielokrotnością identycznej liczby innych wielkości, to suma pierwszych wielkości jest tą samą wielokrotnością sumy drugich.
Twierdzenie 2.
    Jeśli pierwsza wielkość jest taką wielokrotnością drugiej jak trzecia czwartej, a piąta także jest taką samą wielokrotnością drugiej jak szósta czwartej, to suma pierwszej i piątej jest taką samą wielokrotnością drugiej jak suma trzeciej i szóstej jest wielokrotnością czwartej.
Twierdzenie 3.
    Jeżeli pierwsza wielkość jest taką wielokrotnością drugiej, jaką trzecia jest czwartej i jeżeli wzięte są równe wielokrotności pierwszej i trzeciej, wtedy wzięte wielkości też są równymi wielokrotnościami odpowiednio wielkości drugiej i czwartej.
Twierdzenie 4.
W OPRACOWANIU
Twierdzenie 5.
    Jeśli pierwsza wielkość jest taką wielokrotnością drugiej wielkości jak odjęta część pierwszej wielkości jest wielokrotnością odjętej części drugiej wielkości, to reszta pierwszej też jest tą samą wielokrotnością reszty drugiej.
Twierdzenie 6.
    Jeżeli dwie wielkości są takimi samymi wielokrotnościami dwu danych wielkości i dowolne odjęte części tych wielkości są równie wielokrotne dwóm danym wielkościom, to reszty pierwszych wielkości są bądź równe danym, bądź wielokrotne względem nich.
Twierdzenie 7.
    Równe wielkości mają z inną wielkością taki sam stosunek oraz inna wielkość ma z równymi wielkościami ten sam stosunek.
Twierdzenie 8.
W OPRACOWANIU
Twierdzenie 9.
    Wielkości będące w równym stosunku względem danej wielkości, są równe między sobą oraz wielkości, względem których dana wielkość jest w równym stosunku, są równe między sobą.
Twierdzenie 10.
    Z wielkości będących w stosunku do danej wielkości ta, która ma większy stosunek jest większa; natomiast ta, do której dana wielkość ma większy stosunek, jest mniejsza.
Twierdzenie 11.
    Stosunki wielkości, równe pewnemu stosunkowi, są równe między sobą.
Twierdzenie 12.
W OPRACOWANIU
Twierdzenie 13.
    Jeśli pierwsza wielkość jest w takim samym stosunku do drugiej jak trzecia do czwartej i trzecia jest w większym stosunku ma do czwartej, niż piąta do szóstej, to pierwsza jest w większym stosunku do drugiej, niż piąta do szóstej.
Twierdzenie 14.
    Jeśli pierwsza wielkość jest w takim samym stosunku do drugiej jak trzecia do czwartej, i pierwsza jest większa od trzeciej, to druga jest również większa od czwartej; jeżeli jest równa, to równa; jeżeli mniejsza, to mniejsza.
Twierdzenie 15.
    Części wielkości są w tym samym stosunku, co ich równe wielokrotności.
Twierdzenie 16.
W OPRACOWANIU
Twierdzenie 17.
    Jeśli wielkości złożone są proporcjonalne, to także wielkości rozdzielone są proporcjonalne.
Twierdzenie 18.
    Jeśli wielkości rozdzielone są proporcjonalne, to złożone są także proporcjonalne.
Twierdzenie 19.
    Jeżeli całość ma się do całości jak jej część do części, to reszta także ma się do reszty jak całość do całości. Wniosek
    Jeśli wielkości wzięte razem są proporcjonalne, to są też proporcjonalne w konwersji.
Twierdzenie 20.
W OPRACOWANIU
Twierdzenie 21.
    Jeżeli dane są trzy wielkości i inne, równe ich wielokrotności, które wzięte po dwie są w tym samym stosunku oraz proporcja ich jest mieszana, to wtedy, jeśli pierwsza wielkość jest większa od trzeciej, to czwarta jest również większa od szóstej; jeśli równa, to równa; jeśli mniejsza, to mniejsza.
Twierdzenie 22.
W OPRACOWANIU
Twierdzenie 23.
    Jeżeli dane są trzy wielkości i inne, równe ich wielokrotności, które wzięte po dwie są w tym samym stosunku oraz proporcja ich jest mieszana, to są wtedy w tym samym stosunku ex aeguali.
Twierdzenie 24.
W OPRACOWANIU
Twierdzenie 25.
    Jeśli cztery wielkości są proporcjonalne, to suma największej i ostatniej wielkości jest większa niż suma pozostałych dwóch. 



\section{Księga szósta, o Iksach Dupiksach}

Definicja 1.
    Figury podobne prostoliniowe to takie, których kąty są odpowiednio równe, a boki przy tych kątach proporcjonalne.
Definicja 2.
    Dwie figury są do siebie odwrotne kiedy boki przy odpowiadających sobie kątach są odwrotnie proporcjonalne.
Definicja 3.
    Mówi się, że linia prosta została podzielona w złoty sposób wtedy, kiedy cała linia ma się do większego odcinka tak jak większy do mniejszego.
Definicja 4.
    Wysokością jakiejkolwiek figury jest prostopadła narysowana od wierzchołka do podstawy.


Księga VI - Twierdzenia:

Twierdzenie 1.
    Trójkąty i równoległoboki, które mają tę samą wysokość, mają się tak jeden do drugiego jak ich podstawy.
Twierdzenie 2.
    Jeśli linia prosta jest narysowana równolegle do jednego z boków trójkąta, wówczas przecina ona boki trójkąta proporcjonalnie; a jeśli boki trójkąta przecięte są proporcjonalnie, wówczas linia łącząca punkty odcinka jest równoległa do pozostałego boku trójkąta.
Twierdzenie 3.
    Jeśli kąt trójkąta jest podzielony na pół linią prostą przecinającą podstawę, wówczas odcinki podstawy maja tę samą proporcję jak pozostałe boki trójkąta; a jeśli odcinki podstawy mają taką samą proporcję jak pozostałe boki trójkąta, wówczas linia prosta łącząca wierzchołek z punktem odcinka dzieli kąt trójkąta na pół.
Twierdzenie 4.
    W trójkątach równokątnych boki przy kątach równych są proporcjonalne gdzie odpowiadające sobie boki leżą na przeciw kątów równych.
Twierdzenie 5.
    Jeśli dwa trójkąty mają boki proporcjonalne wtedy są trójkątami równokątnymi z kątami równymi leżącymi naprzeciw właściwych boków.
Twierdzenie 6.
    Jeśli dwa trójkąty mają jeden kąt równy drugiemu kątowi i boki przy kątach równych proporcjonalne, wówczas trójkąty te są równokątne i mają te kąty równe naprzeciw odpowiadających boków.
Twierdzenie 7.
    Jeżeli dwa trójkąty mają jeden kąt równy drugiemu kątowi, boki przy innych kątach proporcjonalne, a pozostałe kąty albo mniejsze lub nie mniejsze niż kąt prosty, wtedy trójkąty są równokątne i mają kąty równe przy bokach, które są proporcjonalne.
Twierdzenie 8.
    Jeśli w trójkącie prostokątnym narysowana jest prostopadła od kąta prostego do podstawy, wówczas trójkąty przyległe do prostopadłej są podobne zarówno do całości jak i do siebie nawzajem.
    Wniosek:
    Jeśli w trójkącie prostokątnym narysowana jest prostopadła od kąta prostego do podstawy, wówczas linia prosta tak narysowana tworzy proporcję pomiędzy częściami podstawy.
Twierdzenie 9.
    Wyciąć przepisaną część z danej linii prostej.
Twierdzenie 10.
    Wyciąć daną nie wyciętą linię prostą podobną do danej wyciętej linii prostej.
Twierdzenie 11.
    Znaleźć trzecią proporcjonalną do dwóch danych linii prostych.
Twierdzenie 12.
    Znaleźć czwartą proporcjonalną do trzech danych linii prostych.
Twierdzenie 13.
    Znaleźć najlepszą (złotą) proporcję dwóch danych linii prostych.
Twierdzenie 14.
    W równych i równokątnych równoległobokach boki przy równych kątach są odwrotnie proporcjonalne, a równokątne równoległoboki w których boki przy kątach równych są odwrotnie proporcjonalne są równe.
Twierdzenie 15.
    W trójkątach równych które mają jeden kąt równy drugiemu kątowi boki przy katach równych są odwrotnie proporcjonalne; a te trójkąty, które mają jeden kąt równy drugiemu kątowi, i w których boki przy katach równych są odwrotnie proporcjonalne, są równe.
Twierdzenie 16.
    W OPRACOWANIU
Twierdzenie 17.
    W OPRACOWANIU
Twierdzenie 18.
    Aby opisać prostoliniową figurę podobną i podobnie położoną do danej figury prostoliniowej na danej linii prostej.
Twierdzenie 19.
    Trójkąty podobne odpowiadają sobie w zduplikowanym podziale odpowiednich boków.
    Wniosek:
    Jeśli trzy linie proste są proporcjonalne, wówczas pierwsza ma się do trzeciej jak figura opisana na pierwszej ma się do tego co jest podobne i podobnie opisane na drugiej.
Twierdzenie 20.
    W OPRACOWANIU
    Wniosek:
    W OPRACOWANIU
Twierdzenie 21.
    Figury, które są podobne do tej samej figury prostoliniowej są także podobne do siebie nawzajem.
Twierdzenie 22.
    Jeśli cztery linie proste są proporcjonalne, wówczas figury prostoliniowe podobne i podobnie opisane na nich są proporcjonalne; i jeśli figury prostoliniowe podobne i podobnie opisane na nich są proporcjonalne, wówczas same linie proste są także proporcjonalne.
Twierdzenie 23.
    W OPRACOWANIU
Twierdzenie 24.
    W jakimkolwiek równoległoboku równoległoboki przy średnicy są podobne zarówno do całości jak i do siebie nawzajem.
Twierdzenie 25.
    Skonstruować figurę podobną do danej figury prostoliniowej i równej do drugiej.
Twierdzenie 26.
    Jeśli z równoległoboku wzięty jest równoległobok podobny i podobnie położony do całości i mający wspólny kąt z nim, wówczas ma tę samą średnicę z całością.
Twierdzenie 27.
    W OPRACOWANIU
Twierdzenie 28.
    Umieścić równoległobok równy danej figurze prostoliniowej na linii prostej ale zmniejszony przez równoległobok podobny do danego; dlatego też dana figura prostoliniowa nie może być większa niż równoległobok opisany na połowie linii prostej i podobny do danego równoległoboku.
Twierdzenie 29.
    Umieścić równoległobok równy danej figurze prostoliniowej na danej linii prostej ale powiększając ją o równoległobok podobny do danego.
Twierdzenie 30.
    Podzielić daną skończoną linię w złoty sposób.
Twierdzenie 31.
    W trójkątach prostokątnych figura utworzona na boku leżącym na przeciw kąta prostego równa się sumie podobnych i podobnie położonych figur na bokach zawierających kąt prosty.
Twierdzenie 32.
    Jeśli dwa trójkąty mające dwa boki proporcjonalne do dwóch boków są umieszczone razem w jednym kącie tak że ich odpowiadające sobie boki są także równoległe, wówczas pozostałe boki trójkątów są w linii prostej.
Twierdzenie 33.
    Kąty w równych okręgach mają taką samą proporcję jak obwody kół na których są położone bez względu na to czy leżą w środku czy na obwodzie.
Twierdzenie 34.
    W OPRACOWANIU
Twierdzenie 35.
    W OPRACOWANIU 



\section{Księga dziesiąta, o Iksach Dupiksach}

    Definicja 1.
        Te długości, które da się zmierzyć tą samą miarą zwane są współmierne a te, które nie mogą mieć wspólnej miary zwane są niewspółmierne.
    Definicja 2.
        Linie proste są współmierne w kwadracie, kiedy kwadraty są mierzone na nich tym samym obszarem, a niewspółmierne w kwadracie kiedy kwadraty na nich nie mogą mieć żadnego obszaru jako wspólną miarę.
    Definicja 3.
        Na podstawie tego przypuszczenia udowadniamy, że istnieją nieograniczone sumy odcinków, które są współmierne i niewspółmierne, w tej kolejności niektóre tylko w długości a inne w kwadracie, z przydzielonym odcinkiem. Niech ten przydzielony odcinek będzie wymierny, i te odcinki, które są z nim współmierne, czy w długości i kwadracie, lub tylko w kwadracie, także zwane wymiernymi, a te odcinki, które są z nim niewspółmierne są zwane niewymierne.
    Definicja 4.
        Niech kwadrat na tej przydzielonej linii będzie nazwany wymiernym, a te obszary, które są z nim współmierne będą także zwane wymierne, ale te obszary, które są z nim niewspółmierne będą niewymierne, a linie proste, które je tworzą także są niewymierne, to znaczy w przypadku w którym obszary są kwadratami to boki, ale w przypadku w jakim one są jakimikolwiek innymi figurami to proste linie na których opisane są kwadratami równymi do nich.
    
    
    Księga X - Twierdzenia I:
    
    Twierdzenie 1.
        Dwie nierówne wielkości będąc ustalone, jeśli od większej odjęta jest wielkość większa za od jej połowy i od tego co zostaje także odjęta wielkość od jej połowy, i proces ten będzie powtarzany to pozostanie wielkość mniejsza od tej mniejszej ustalonej wielkości. To twierdzenie może być podobnie udowodnione, nawet kiedy części odejmowane to tylko połówki.
    Twierdzenie 2.
        Jeżeli, kiedy mniejsza z dwóch nierównych wielkości jest wciąż odejmowana od tej większej wielkości, i to co zostaje nigdy nie mierzy to poprzednie, to te dwie wielkości są wtedy niewspółmierne.
    Twierdzenie 3.
        Żeby znaleźć największą wspólną miarę dwóch danych odległości.
        Wniosek:
        Jeżeli odległość mierzy dwie odległości to ona także mierzy ich największą wspólna miarę.
    Twierdzenie 4.
        Żeby znaleźć największa wspólną miarę trzech danych odległości.
        Wniosek:
        Jeżeli odległość mierzy trzy odległości to ona także mierzy ich największą wspólna miarę. Największa wspólna miara może być podobnie znaleziona dla większych ilości odległości, i wniosek przedłużony.
    Twierdzenie 5.
        Współmierne wielkości mają wobec siebie stosunek taki sam jak liczba do liczby.
    Twierdzenie 6.
        Jeżeli dwie wielkości maja wobec siebie stosunek jak liczba do liczby to te wielkości są niewspółmierne.
    Twierdzenie 7.
        Niewspółmierne wielkości nie mają wobec siebie stosunku, który liczba ma do liczby.
    Twierdzenie 8.
        Jeżeli wielkości nie mają wobec siebie stosunku, który ma liczba wobec liczby to te wielkości są niewspółmierne.
    Twierdzenie 9.
        Kwadraty na liniach prostych współmierne w długości maja wobec siebie stosunek, który maja dwie liczby do kwadratu wobec siebie, i kwadraty, które mają wobec siebie stosunek jak dwie liczby do kwadratu wobec siebie, także maja boki współmierne w długości. Ale kwadraty na liniach prostych niewspółmiernych w długości nie maja wobec siebie stosunku jak dwie liczby do kwadratu wobec siebie i kwadraty, które nie maja do siebie stosunku jak dwie liczby do kwadratu wobec siebie, także nie maja boków współmiernych w długości.
        Wniosek:
        Linie proste współmierne w wielkości są także współmierne w kwadracie, ale te współmierne w kwadracie nie zawsze są współmierne w długości.
        Lemat:
        Liczby podobne płaszczyznowo maja do siebie stosunek jak dwie liczby do kwadratu wobec siebie. Jeśli dwie liczby wobec siebie maja stosunek jak dwie liczby do kwadratu wobec siebie to są liczby podobnie płaszczyznowo.
        Następstwo 2:
        Liczby, które nie są liczbami podobnymi płaszczyznowo to znaczy liczby, które nie mają proporcjonalnych boków, nie maja do siebie stosunku jak dwie liczby do kwadratu wobec siebie.
    Twierdzenie 10.
        Żeby znaleźć dwie linie proste niewspółmierne, jedna tylko w długości, a druga tylko w kwadracie z przydzieloną linia prostą.
    Twierdzenie 11.
        Jeśli cztery z wielkości są proporcjonalne, i pierwsza jest współmierna z drugą, to trzecia jest także współmierna z czwartą, ale jeżeli pierwsza jest niewspółmierna z drugą, to trzecia jest także niewspółmierna z czwartą.
    Twierdzenie 12.
        Wielkości współmierne z ta samą wielkością są ze sobą współmierne.
    Twierdzenie 13.
        Jeżeli dwie wielkości są współmierne i jedna z nich jest niewspółmierna z jakąkolwiek wielkością, to pozostająca jest także niewspółmierna z tym samym.
    Twierdzenie 14.
        Jeżeli cztery linie są proporcjonalne i kwadrat jest na pierwszej jest większy od kwadratu na drugiej, o kwadrat na linii prostej współmiernej z pierwszą to kwadrat na trzeciej jest także większy niż kwadrat na czwartej, o kwadrat na trzeciej linii współmiernej z trzecim, i jeśli kwadrat na pierwszej linii jest większy niż kwadrat na drugiej linii, o kwadrat na linii prostej, niewspółmiernej z pierwszą, to kwadrat na trzeciej linii jest także większy od kwadratu na czwartej linii, o kwadrat na trzeciej linii niewspółmiernej z trzecią.
        Lemat:
        Mając dwie nierówne linie proste, aby znaleźć o jaki kwadrat , kwadrat na większej linii jest większy od kwadratu na mniejszej. Mając dwie linie proste, aby znaleźć linię prostą na, której kwadrat równa się suma kwadratów na nich.
    Twierdzenie 15.
        Jeżeli współmierne wielkości są dodane razem, to całość jest także współmierna z każdą z nich. A, jeżeli całość jest współmierna z jedną z nich, to oryginalne wielkości są również współmierne.
    Twierdzenie 16.
        Jeżeli dwie niewspółmierne wielkości są dodane razem ,to suma jest także niewspółmierna z każdą z nich, ale jeżeli suma jest niewspółmierna z jedną z nich, te oryginalne wielkość są także niewspółmierne.
    Twierdzenie 17.
        Jeśli są dwie nierówne linie proste, i do większej jest dodany równoległobok, równy do czwartej części kwadratu na mniejszej, ale brakującej kwadratu, jeżeli podzieli się to na części współmierne w długości to kwadrat na większej jest większy od kwadratu na mniejszej o kwadrat na linii prostej współmiernej z mniejszą. Jeżeli kwadrat na większej jest większy od kwadratu na mniejszej o kwadrat na linii prostej współmiernej z mniejszą. Jeśli jest dodany do większej równoległobok równy do czwartej części na mniejszej brakującej o kwadrat, to jest podzielony na części współmierne w długości.
        Lemat:
        Jeżeli do jakiejkolwiek linii prostej jest dodany równoległobok, któremu brakuje kwadratu, to ten dodany równoległobok jest równy do prostokąta mieszczącemu się w odcinkach linii prostych wynikłych z dodawania.
    Twierdzenie 18.
        Jeżeli są dwie nierówne linie proste, i do większej dodamy równoległobok równy do czwartej części kwadratu na mniejszej, której brakuje kwadrat. I jeżeli to się dzieli na niewspółmierne części, to kwadrat na większej jest większy od kwadratu na mniejszej o kwadrat na linii prostej niewspółmiernej z większą. I jeśli kwadrat na większej jest większy od kwadratu na mniejszej o linię prostą niewspółmiernej z większą i jeśli dodamy do większej równoległobok równy do czwartej części kwadratu na mniejszej, ale brakującej o jeden kwadrat to podzieli się to na niewspółmierne części.
    Twierdzenie 19.
        Prostokąt mieszczący się w wymiernych liniach prostych współmiernych w długości jest wymierny.
    Twierdzenie 20.
        Jeżeli wymierny obszar jest dodany do wymiernego odcinka to jako szerokość da to odcinek wymierny i współmierny w długości z odcinkiem do, którego jest dodany.
    Twierdzenie 21.
        Prostokąt mieszczący się w wymiernych odcinkach współmiernych tylko w kwadracie, jest niewymierny i boki kwadratu do niego równego są niewymierne. Niech one nazywają się środkowe (średnie).
    Twierdzenie 22.
        Kwadrat na środkowym odcinku, jeśli dodamy do wymiernego odcinka daje jako szerokość odcinek wymierny i niewspółmierny w długości z tym do, którego jest dodany.
        Lemat:
        Jeżeli są dwa odcinki, to pierwszy jest do drugiego tak, jak kwadrat na pierwszym do prostokąta mieszczącego się między tymi odcinkami.
    Twierdzenie 23.
        Odcinek współmierny ze środkowym odcinkiem jest środkowy.
        Wniosek:
        Obszar współmierny ze środkowym obszarem jest środkowy.
    Twierdzenie 24.
        Prostokąt mieszczący się w środkowych odcinkach wymiernych w długości jest środkowy.
    Twierdzenie 25.
        Prostokąt mieszczący się w środkowym odcinku współmiernym w kwadracie jest albo wymierny albo środkowy.
    Twierdzenie 26.
        Środkowy obszar nie przekracza środkowego obszaru o wymierny obszar. 
    Twierdzenie 27.
        Aby znaleźć środkowe odcinki współmierne w kwadracie tylko, w których mieści się prostokąt środkowy.
    Twierdzenie 28.
        Aby znaleźć środkowe odcinki współmierne w kwadracie tylko, w których mieści się prostokąt.
    Twierdzenie 29.
        Aby znaleźć dwa wymierne odcinki współmierne tylko w kwadracie, tak aby kwadrat na większej był większy od kwadratu na mniejszej o kwadrat na odcinku współmiernym na odcinku z większym.
        Lemat 1:
        Aby znaleźć dwie liczby kwadratowe i aby ich suma także była kwadratowa.
        Lemat 2:
        Aby znaleźć dwie liczby kwadratowe i aby ich suma nie była kwadratowa.
    Twierdzenie 30.
        Aby znaleźć dwa wymierne odcinki współmierne tylko w kwadracie tak aby kwadrat był większy niż kwadrat na mniejszym o kwadrat na odcinku niewspółmiernego długości z większych.
    Twierdzenie 31.
        Aby znaleźć dwa środkowe odcinki współmierne tylko w kwadracie, w których mieści się środkowy prostokąt, tak aby kwadrat na większym był większy od kwadratu na mniejszym o kwadrat na odcinku współmiernym w długości z większym.
    Twierdzenie 32.
        Aby znaleźć dwa środkowe odcinki współmierne tylko w kwadracie, w którym mieści się środkowy trójkąt, tak aby kwadrat na większej, był większy od kwadratu na mniejszej o kwadrat na odcinku współmierna z większym odcinkiem.
    Twierdzenie 33.
        Aby znaleźć dwa odcinki niewspółmierne w kwadracie, których suma kwadratów na nich była wymierne , ale prostokąt, który się między nimi mieści był środkowy.
    Twierdzenie 34.
        Aby znaleźć dwa odcinki niewspółmierne w kwadracie, w których suma kwadratów na nich była środkowa, ale prostokąt, który się miedzy nimi mieści był wymierny.
    Twierdzenie 35.
        Aby znaleźć dwa odcinki niewspółmierne w kwadracie, gdzie suma kwadratów na nich była środkowa a prostokąt mieszczący się między nimi był środkowy i niewspółmierny z sumą kwadratów na nich.
    Twierdzenie 36.
        Jeżeli dwa wymierne odcinki współmierne tylko w kwadracie są dodane, a całość jest niewymierna, niech będzie wiec nazwana dwumianem.
    Twierdzenie 37.
        Jeżeli dwa wymierne odcinki współmierne tylko w kwadracie, miedzy którymi mieści się środkowy prostokąt, są dodane, i całość jest niewymierna, to niech będzie więc nazwana pierwszym dwumianem odcinka.
    Twierdzenie 38.
        Jeżeli dwa środkowe odcinki współmierne tylko w kwadracie, miedzy którymi mieści się środkowy prostokąt, są dodane, i całość jest niewymierna, to niech będzie więc nazwana drugim dwumianem odcinka.
    Twierdzenie 39.
        Jeżeli dwa odcinki niewspółmierne w kwadracie, gdzie suma kwadratów na nich jest wymierna, ale prostokąt mieszczący sie miedzy nimi jest środkowy, są dodane, to cały odcinek jest niewymierny, niech będzie nazwany wielki.
    Twierdzenie 40.
        Jeśli dwa odcinki niewspółmierne w kwadracie, gdzie suma kwadratów na nich jest środkowa, ale prostokąt mieszczący się między nimi jest wymierny, i są dodane, to cały odcinek nazwijmy bokiem wymiernego plus środkowy obszar.
    Twierdzenie 41.
        Jeżeli dwa odcinki niewymierne w kwadracie, gdzie suma kwadratów na nich jest środkowa prostokąt mieszczący się między nimi jest środkowy i także niewymierny z suma kwadratów na nich, są dodane, to cały odcinek jest niewymierny. Niech będzie wiec nazywany bok sumy dwóch środkowych obszarów.
    Twierdzenie 42.
        Dwumienny odcinek jest podzielony na swoje terms tylko w jednym punkcie.
    Twierdzenie 43.
        Pierwszy dwu środkowy odcinek jest podzielony w jednym i tym samym miejscu.
    Twierdzenie 44.
        Drugi dwu środkowy odcinek, jest dzielony tylko w jednym miejscu.
    Twierdzenie 45.
        Wielki segment jest dzielony tylko w jednym miejscu.
    Twierdzenie 46.
        Bok wymiernego dodany do środkowego obszaru jest dzielony tylko w jednym miejscu.
    Twierdzenie 47.
        Bok sumy dwóch środkowych obszarów jest dzielony tylko w jednym miejscu.
    
    
    Księga X - Definicje II:
    
    Definicja 5.
        Danym wymierny odcinek i dwumienną, podzielone na terms, tak aby kwadrat na większym term był większy niż kwadrat na mniejszym, o kwadrat na odcinku współmiernym w długości z większym, wtedy jeśli większy term jest współmierny w długości z danym wcześniej wymiernym odcinkiem. Niech całość będzie nazwana pierwszym dwumiennym odcinkiem.
    Definicja 6.
        Ale, mniejszy term jest współmierny w długości z wymiernym odcinkiem, i niech całość będzie drugim dwumiennym.
    Definicja 7.
        Jeśli żaden z term jest współmierny w długości z danym na początku wymiernym odcinkiem, to niech całość będzie zwana trzecim dwumiennym.
    Definicja 8.
        Jeśli kwadrat na większym term jest większy niż kwadrat na mniejszym, o kwadrat na odcinku niewspółmiernym w długości z większym, i jeśli większy term jest współmierny w długości z danym na początku wymiernym odcinkiem, to niech całość będzie zwana czwartym dwumiennym.
    Definicja 9.
        A jeśli mniejszy to będzie zwane piątym dwumiennym.
    Definicja 10.
        A jeśli żaden, to będzie zwana szóstym dwumiennym.
    
    
    Księga X - Twierdzenia II:
    
    Twierdzenie 48.
        Znalezienie pierwszej dwumiennej linii.
    Twierdzenie 49.
        Znalezienie drugiej dwumiennej linii.
    Twierdzenie 50.
        Znalezienie trzeciej dwumiennej linii.
    Twierdzenie 51.
        Znalezienie czwartej dwumiennej linii.
    Twierdzenie 52.
        Znalezienie piątej dwumiennej dwumiennej linii.
    Twierdzenie 53.
        Znalezienie szóstej dwumiennej linii.
    Twierdzenie 54.
        Jeśli obszar mieści się w wymiernym odcinku i pierwszej dwumiennej, wtedy bok tego obszaru jest niewymierny odcinkiem, który jest dwu środkowy.
    Twierdzenie 55.
        Jeśli obszar mieści się w wymiernym odcinku i drugiej dwumiennej, wtedy bok tego obszaru jest niewymierny odcinkiem, który jest pierwszą dwu środkową.
    Twierdzenie 56.
        Jeśli obszar mieści się w wymiernym odcinku i trzeciej dwumiennej, wtedy bok tego obszaru jest niewymierny z odcinkiem, który jest drugą dwu środkową.
    Twierdzenie 57.
        Jeśli obszar mieści się w wymiernym odcinku i czwartej dwumiennej, wtedy bok tego obszaru jest niewymierny z odcinkiem, który jest wielki.
    Twierdzenie 58.
        Jeśli obszar mieści się w wymiernym odcinku i piątej dwumiennej, wtedy bok tego obszaru jest niewymierny z odcinkiem, który jest bokiem wymiernym dodanym do środkowego obszaru.
    Twierdzenie 59.
        Jeśli obszar mieści się w wymiernym odcinku i szóstej dwumiennej, wtedy bok tego obszaru jest niewymierny z odcinkiem, który jest bokiem sumy dwóch środkowych obszarów.
    Twierdzenie 60.
        Kwadrat na dwumiennym odcinku dodany do wymiernego odcinka daje jako szerokość pierwszy dwumienny.
        Lemat:
        Jeśli odcinek jest przecięte w nierówne części to suma kwadratów na tych nierównych częściach jest większa niż dwa razy prostokąt mieszczący się w nierównych częściach.
    Twierdzenie 61.
        Kwadrat na pierwszym dwumiennym odcinku dodany do wymiernego odcinka daje jako szerokość drugi dwumienny.
    Twierdzenie 62.
        Kwadrat na drugim dwumiennym odcinku dodawany do wymiernego odcinka daje jako szerokość trzeci dwumienny.
    Twierdzenie 63.
        Kwadrat na wielkim odcinku dodawany do wymiernego odcinka daje szerokość jako czwarty dwumienny.
    Twierdzenie 64.
        Kwadrat na boku wymiernej, dodanej do środkowego obszaru, dodane do wymiernego odcinka daje, jako szerokość piąty dwumienny.
    Twierdzenie 65.
        Kwadrat na boku sumy dwóch środkowych obszarów, dodany do wymiernego odcinka daje jako szerokość szósty dwumienny.
    Twierdzenie 66.
        Odcinek współmierny z dwumiennym odcinkiem jest sam w sobie dwumienny i w takiej samej kolejności.
    Twierdzenie 67.
        Odcinek współmierny z dwu środkowym odcinkiem jest sam w sobie dwu środkowy i w tej samej kolejności.
    Twierdzenie 68.
        Odcinek współmierny z wielkim odcinkiem jest sam w sobie także wielki.
    Twierdzenie 69.
        Odcinek współmierny z bokiem wymiernym dodany do środkowego obszaru jest sam w sobie także bokiem wymiernym dodanym do środkowego obszaru.
    Twierdzenie 70.
        Odcinek współmierny z bokiem sumy dwóch środkowych obszarów jest bokiem sumy dwóch środkowych obszarów.
    Twierdzenie 71.
        Jeśli wymierny i środkowy są dodane to cztery niewymierne odcinki powstają, a dokładniej to dwumienny, albo pierwszy dwu środkowy, albo wielki, albo bok wymiernego dodanego do środkowego obszaru.
    Twierdzenie 72.
        Jeżeli dwa środkowe obszary niewspółmierne ze sobą są dodane, to pozostające dwa niewymierne odcinki powstają, a są to, albo dwu środkowe, albo bok sumy dwóch środkowych obszarów.
    Twierdzenie 73.
        Jeśli od wymiernego odcinka odejmiemy wymierny odcinek współmierny z całością tylko w kwadracie to reszta jest niewymierna, więc niech nazwana będzie apotome.
    Twierdzenie 74.
        Jeżeli od środkowego odcinka jest odjęty środkowy odcinek współmierny z całością, tylko w kwadracie i w którym mieści się z całością, wymierny prostokąt, to reszta jest niewymierna, i niech będzie nazwana pierwsze apotome środkowego odcinka.
    Twierdzenie 75.
        Jeśli od środkowego odcinka jest odjęty środkowy odcinek, który jest współmierny z całością tylko w kwadracie, i w którym mieści się z całością środkowy prostokąt, to reszta będzie niewymierna i nazwana drugie apotome środkowego odcinka. 
    Twierdzenie 76.
        Jeśli od odcinka jest odjęty odcinek który jest niewspółmierny w kwadracie z całością i w który z całością tworzy sumę kwadratu na nich dodanych razem wymiernych ale prostokąt mieszczący się w nich jest środkowy to reszta jest niewymierna i niech będzie zwana mniejszą.
    Twierdzenie 77.
        Jeśli od odcinka odjęty jest odcinek, który jest niewspółmierny w kwadracie z całością, i który z całością tworzy sumę kwadratów na nich środkową, ale dwa razy prostokąt mieszczący się w nich wymierny, reszta jest niewymierna, i niech będzie nazwana tym co daje wymiernym obszarem środkową całość.
    Twierdzenie 78.
        Jeśli od odcinka jest odjęty niewspółmierny w kwadracie z całością, i który z całością tworzy sumę kwadratów na nich środkową, dwa razy prostokąt, mieszczący się miedzy nimi środkowy i także, kwadraty na nich niewspółmierne z dwa razy prostokątem mieszczącym się miedzy nimi, to reszta jest niewymierna i niech będzie nazwana to, które daje ze środkowym obszarem środkowa całość.
    Twierdzenie 79.
        Do apotome tylko jeden wymierny odcinek może być załączony , który jest współmierny z całością tylko w kwadracie.
    Twierdzenie 80.
        Do pierwszego apotoma środkowego odcinka, tylko jeden środkowy odcinek może być załączony, który jest współmierny z całością tylko w kwadracie, i w którym mieści się z całością środkowy prostokąt.
    Twierdzenie 81.
        Do drugiego apotome środkowego odcinka, tylko jeden środkowy odcinek może być załączony, który jest współmierny z całością, tylko w kwadracie, i który mieści z całością środkowy prostokąt.
    Twierdzenie 82.
        Do mniejszego odcinka tylko jeden odcinek może być załączony, który jest niewspółmierny w kwadracie z całością, i który tworzy z całością sumę kwadratów na nich wymierną, ale dwa razy tyle co prostokąt mieszczący się miedzy nimi, środkowy.
    Twierdzenie 83.
        Do odcinka, który tworzy z wymiernym obszarem środkową całość, tylko jeden odcinek może być załączony, który jest niewspółmierny kwadracie z całym odcinkiem, i który z całym odcinkiem tworzy sumę kwadratów na nich środkową, a dwa razy prostokąt mieszczący się miedzy nimi wymierny.
    Twierdzenie 84.
        Do odcinka, który tworzy ze środkowym obszarem środkową całość, tylko jeden odcinek może być dołączony, który jest niewspółmierny w kwadracie z całym odcinkiem, i który z całym odcinkiem tworzy sumę kwadratów na nich środkowa i dwa razy prostokąt mieszczący się w nich środkowy i niewspółmierny z suma kwadratów na nich.
    
    
    Księga X - Definicje III:
    
    Definicja 11.
        Dany wymierny odcinek i apotome, jeśli kwadrat na całości jest większy, niż kwadrat na załączeniu o kwadrat na odcinku współmiernym w długości z całością i całość jest współmierna w długości z wymiernym odcinkiem danym na początku, niech apotome będzie zwane pierwsze apotome.
    Definicja 12.
        Ale, jeśli załączenie jest współmierne z wymiernym odcinkiem danym na początku i kwadrat na całości jest większy niż kwadrat na załączeniu, o kwadrat na odcinku współmiernym z całością niech to apotome będzie zwane drugie apotome.
    Definicja 13.
        Ale, jeśli żaden nie jest współmierny w długości z wymiernym odcinkiem danym na początku i kwadrat na całości jest większy, niż kwadrat na załączeniu,, o kwadrat na odcinku współmiernym z całością niech to apotome, będzie się nazywało trzecim apotome.
    Definicja 14.
        Jeżeli kwadrat na całości jest większy, niż kwadrat na załączeniu o kwadrat na odcinku niewspółmiernym z całością, to jeśli całość jest współmierna w długości z wymiernym odcinkiem danym na początku to, niech apotome będzie zwane czwartym apotome.
    Definicja 15.
        Jeśli załączenie jest współmierne, to będzie to piąte apotome. 
    Definicja 16.
        Ale, jeśli żadne to będzie to szóste apotome.
    
    
    Księga X - Twierdzenia III:
    
    Twierdzenie 85.
        Znalezienie pierwszego apotome.
    Twierdzenie 86.
        Znalezienie drugiego apotome.
    Twierdzenie 87.
        Znalezienie trzeciego apotome.
    Twierdzenie 88.
        Znalezienie czwartego apotome.
    Twierdzenie 89.
        Znalezienie piątego apotome.
    Twierdzenie 90.
        Znalezienie szóstego apotome.
    Twierdzenie 91.
        Jeżeli obszar jest umieszczony na wymiernym odcinku i pierwszej apotome, wtedy bok obszaru jest apotome.
    Twierdzenie 92.
        Jeżeli obszar jest umieszczony na wymiernym odcinku i drugiej apotome, wtedy bok obszaru będzie pierwszą apotome i środkowego odcinka.
    Twierdzenie 93.
        Jeżeli obszar jest umieszczony na wymiernym odcinku i trzecim apotome to bok tego obszaru jest drugim apotome środkowego odcinka.
    Twierdzenie 94.
        Jeżeli obszar jest umieszczony na wymiernym odcinku i czwartym apotome to bok tego obszaru jest mniejszy. 
    Twierdzenie 95.
        Jeżeli obszar jest umieszczony na wymiernym odcinku i piąty apotome to bok tego obszaru jest odcinkiem, który daje z wymiernym obszarem środkowa całość.
    Twierdzenie 96.
        Jeżeli obszar jest umieszczony na wymiernym odcinku i szóstej apotome to bok tego obszaru jest odcinkiem, który daje ze środkowym obszarem środkowa całość.
    Twierdzenie 97.
        Kwadrat na apotome, środkowego odcinka dodany do wymiernego odcinka, daje jako szerokość pierwsza apotome.
    Twierdzenie 98.
        Kwadrat na pierwszej apotome, środkowego odcinka dodany do wymiernego odcinka daje, jako szerokość drugie apotome.
    Twierdzenie 99.
        Kwadrat na drugiej apotome, środkowego odcinka dodany do wymiernego odcinka daje jako szerokość trzecie apotome.
    Twierdzenie 100.
        Kwadrat na mniejszym odcinku, dodany do wymiernego odcinka daje, jako szerokość czwarte apotome.
    Twierdzenie 101.
        Kwadrat na odcinku, który daje z wymiernym obszarem środkową całość, jeśli dodany do wymiernego odcinka daje, jako szerokość piąte apotome.
    Twierdzenie 102.
        Kwadrat na odcinku, który daje ze środkowym obszarem środkową całość, jeśli dodany do wymiernego odcinka daje jako szerokość szóste apotome.
    Twierdzenie 103.
        Odcinek współmierny w długości z apotome jest apotome i w tej samej kolejności.
    Twierdzenie 104.
        Odcinek współmierny z apotome środkowego odcinka jest apotome środkowego odcinka i w tej samej kolejności.
    Twierdzenie 105.
        Odcinek współmierny z mniejszym odcinkiem jest mniejszy.
    Twierdzenie 106.
        Odcinek współmierny z tym co daje z wymiernym obszarem środkową całość jest odcinkiem, który daje z wymiernym obszarem środkową całość.
    Twierdzenie 107.
        Odcinek współmierny z tym co tworzy ze środkowym obszarem, środkową całość jest sama w sobie także odcinkiem, który daje ze środkowym obszarem, środkową całość.
    Twierdzenie 108.
        Jeśli od wymiernego obszaru odjęty jest środkowy obszar, bok pozostającego obszaru staje się jedną z dwóch odcinków, albo apotome, albo odcinek mniejszy.
    Twierdzenie 109.
        Jeżeli od środkowego obszaru odjęty jest wymierny obszar to powstają dwa niewymierne odcinki, albo pierwsze apotome środkowego odcinka, albo odcinek który daje z wymiernym obszarem środkowa całość.
    Twierdzenie 110.
        Jeżeli od środkowego obszaru odjęty jest środkowy obszar niewspółmierny z całością, to dwa pozostające niewymierne odcinki pojawiają się, albo drugie apotome środkowego odcinka, albo odcinek, który daje ze środkowym obszarem środkowa całość
    Twierdzenie 111.
        Apotome nie jest tym samym co dwumienny odcinek.
    Twierdzenie 112.
        Kwadrat na wymiernym odcinku dodany do dwumiennego odcinka daje jako szerokość apotome których terms są niewspółmierne z terms dwumiennego odcinka i także w tym samym stosunku. Także aptome w ten sposób powstające ma ta samą kolejność co dwumienny odcinek.
    Twierdzenie 113.
        Kwadrat na wymiernym odcinku jeśli dodany do apotome tworzy, jako szerokość dwumienny odcinek, który w terms są współmierne z terms z apotome i w tym samym stosunku także dwumienny w ten sposób powstający ma tą samą kolejność co apotome.
    Twierdzenie 114.
        Jeżeli obszar mieszczący się w apotome i dwumiennym odcinku którego terms są współmierne z terms i apotome i w tym samym stosunku to bok tego obszaru jest wymierny.
        Wniosek:
        Możliwe jest aby wymierny obszar mieścił się w niewymiernych odcinkach.
    Twierdzenie 115.
        Ze środkowego odcinka tworzone są niewymierne odcinki z nieograniczona ilością i żaden z nich nie jest taki sam jak następny. 


        
\section{Księga jedenasta, o Iksach Dupiksach}

Definicja 1.
    Bryła to figura, która ma długość, szerokość i głębokość. 
    WYJAŚNIENIE
Definicja 2.
    Bryłę ograniczają płaszczyzny.
    WYJAŚNIENIE
Definicja 3.
    Prosta jest prostopadła do płaszczyzny, gdy tworzy kąty proste z wszystkimi prostymi, które ja przecinają i leżą w danej płaszczyźnie.
    WYJAŚNIENIE
Definicja 4.
    Płaszczyzna jest prostopadła do płaszczyzny, gdy proste poprowadzone w jednej z nich prostopadle do przecięcia płaszczyzn, są prostopadłe do drugiej płaszczyzny.
    WYJAŚNIENIE
Definicja 5.
    Kąt nachylenia prostej do płaszczyzny, to kąt między daną prostą a prostą, która łączy punkt przecięcia danej prostej z płaszczyzną i punkt, który jest przecięciem płaszczyzny i prostej poprowadzonej prostopadle do płaszczyzny przez dowolny punkt należący do danej prostej.
    WYJAŚNIENIE
Definicja 6.
    Mówimy, że nachylenie płaszczyzn jest podobne do nachylenia innych płaszczyzn, gdy ich kąty nachylenia są sobie równe.
    WYJAŚNIENIE
Definicja 7.
    Mówimy, że nachylenie płaszczyzn jest podobne do nachylenia innych płaszczyzn, gdy ich kąty nachylenia są sobie równe.
    WYJAŚNIENIE
Definicja 8.
    Płaszczyzny równoległe to takie, które się nie przecinają.
    WYJAŚNIENIE
Definicja 9.
    Bryły podobne to takie, które są ograniczone taką samą ilością podobnych płaszczyzn.
    WYJAŚNIENIE
Definicja 10.
    Bryły równe i podobne to takie, które ograniczone są taką samą ilością płaszczyzn podobnych i równych co do wielkości.
    WYJAŚNIENIE
Definicja 11.
    Kąt wielościenny jest nachyleniem utworzonym przez więcej niż dwie linie, które się przecinają i nie leżą na tej samej powierzchni, do wszystkich linii, tzn. kąt wielościenny to taki, który jest ograniczony przez więcej niż dwa kąty płaskie, które nie leżą w jednej płaszczyźnie i mają wspólny wierzchołek.
    WYJAŚNIENIE
Definicja 12.
    Ostrosłup jest bryłą ograniczoną przez płaszczyzny, które są poprowadzone z jednej płaszczyzny do jednego punktu.
    WYJAŚNIENIE
Definicja 13.
    Graniastosłup trójkątny jest bryłą ograniczoną przez płaszczyzny, z których dwie, tzn. te przeciwległe, są równe i podobne (przystające) oraz równoległe, podczas gdy pozostałe są równoległobokami.
    WYJAŚNIENIE
Definicja 14.
    Gdy półkole o stałej średnicy obrócimy aż powrócimy do tej samej pozycji, z której zaczęliśmy go obracać, to otrzymaną figurę przyjmujemy jako kulę. 
    WYJAŚNIENIE
Definicja 15.
    Osią kuli jest prosta, która pozostaje stała i, wokół której obracamy półkole.
    WYJAŚNIENIE
Definicja 16.
    Środkiem kuli jest środek półkola
    WYJAŚNIENIE
Definicja 17.
    Średnicą kuli jest dowolny odcinek poprowadzony przez środek i ograniczony po obu stronach przez powierzchnię kuli.
    WYJAŚNIENIE
Definicja 18.
    Jeśli trójkąt prostokątny o stałej jednej z przyprostokątnych obrócimy aż do powrotu do pozycji, z której zaczęliśmy obracać, to otrzymaną figurę nazywamy stożkiem. I, jeśli bok, który pozostał stały, jest równy drugiej przyprostokątnej - tej obracanej, to stożek jest prostokątny; jeśli mniejszy, to rozwartokątny; a jeśli większy, to ostrokątny.
    WYJAŚNIENIE
Definicja 19.
    Osią stożka jest prosta, która zawiera stały bok, wokół którego obracaliśmy trójkąt.
    WYJAŚNIENIE
Definicja 20.
    Podstawą jest koło otrzymane przez bok, który był obracany.
    WYJAŚNIENIE
Definicja 21.
    Jeśli prostokątny równoległobok o jednym stałym boku obrócimy aż powrócimy do pozycji, z której zaczęliśmy go obracać, to otrzymana figura jest walcem.
    WYJAŚNIENIE
Definicja 22.
    Osią walca jest prosta zawierająca stały bok, wokół którego obracaliśmy równoległobok (prostokat).
    WYJAŚNIENIE
Definicja 23.
    Podstawami są koła zakreślone przez dwa przeciwległe boki, które były obracane.
    WYJAŚNIENIE
Definicja 24.
    Podobne stożki i walce to takie, w których osie i średnice podstaw są proporcjonalne.
    WYJAŚNIENIE
Definicja 25.
    Sześcian jest bryłą ograniczoną przez sześć równych kwadratów.
    WYJAŚNIENIE
Definicja 26.
    Ośmiościan jest bryłą ograniczoną przez osiem równych trójkątów równobocznych.
    WYJAŚNIENIE
Definicja 27.
    Dwudziestościan jest bryłą ograniczoną przez dwadzieścia równych trójkątów równobocznych.
    WYJAŚNIENIE
Definicja 28.
    Dwunastościan jest bryłą ograniczoną przez dwanaście równych pięciokątów równobocznych i równokątnych (foremnych).
    WYJAŚNIENIE


Księga XI - Twierdzenia:

Twierdzenie 1.
    Część prostej nie może leżeć w płaszczyźnie odniesienia, a część w płaszczyźnie nachylonej do niej pod pewnym kątem.
    WYJAŚNIENIE
Twierdzenie 2.
    Jeśli dwie proste przecinają się, to leżą one w jednej płaszczyźnie; i każdy trójkąt leży w jednej płaszczyźnie.
    WYJAŚNIENIE
Twierdzenie 3.
    Jeśli dwie płaszczyzny przecinają się, to ich przecięciem jest linia prosta.
    WYJAŚNIENIE
Twierdzenie 4.
    Jeśli prosta jest do dwóch prostych przecinających się prostopadła w punkcie ich przecięcia, to jest prostopadła do płaszczyzny na nich rozpiętej.
    WYJAŚNIENIE
Twierdzenie 5.
    Jeśli prosta jest prostopadła do trzech przecinających się prostych w ich punkcie przecięcia, to proste te leżą w jednej płaszczyźnie.
    WYJAŚNIENIE
Twierdzenie 6.
    Jeśli dwie proste są prostopadłe do tej samej płaszczyzny, to są one równoległe względem siebie.
    WYJAŚNIENIE
Twierdzenie 7.
    Jeżeli dwie proste są równoległe i obierzemy dowolne punkty na każdej z tych prostych, to prosta łącząca te punkty leży w tej samej płaszczyźnie co dane proste równoległe.
    WYJAŚNIENIE
Twierdzenie 8.
    Jeżeli dwie proste są równoległe i jedna z nich jest prostopadła do pewnej płaszczyzny, to pozostała prosta jest również prostopadła do tej samej płaszczyzny.
    WYJAŚNIENIE
Twierdzenie 9.
    Proste, które są równoległe do tej samej prostej, ale nie leżą z nią w jednej płaszczyźnie, są do siebie równoległe.
    WYJAŚNIENIE
Twierdzenie 10.
    Jeżeli dwie przecinające się proste są równoległe do dwóch przecinających się, leżących w innej płaszczyźnie, prostych, to tworzą one równe kąty.
    WYJAŚNIENIE
Twierdzenie 11.
    Poprowadzić prostą prostopadłą do danej płaszczyzny z danego, nie leżącego na niej, punktu.
    WYJAŚNIENIE
Twierdzenie 12.
    Z danego punktu na płaszczyźnie wyprowadzić prostą prostopadłą do tej płaszczyzny.
    WYJAŚNIENIE
Twierdzenie 13.
    Z jednego punktu nie można poprowadzić dwóch różnych prostych prostopadłych do tej samej płaszczyzny po tej samej stronie.
    WYJAŚNIENIE
Twierdzenie 14.
    Płaszczyzny, do których ta sama prosta jest prostopadła, są równoległe.
    WYJAŚNIENIE
Twierdzenie 15.
    Jeśli dwie przecinające się proste są równoległe do dwóch prostych przecinających się w innej płaszczyźnie, to płaszczyzny na nich rozpięte są równoległe.
    WYJAŚNIENIE
Twierdzenie 16.
    Jeśli dwie płaszczyzny równoległe przecina dowolna płaszczyzna, to ich przecięcia są równoległe.
    WYJAŚNIENIE
Twierdzenie 17.
    Jeśli dwie proste przecięte są płaszczyznami równoległymi, to dzielą się one w tym samym stosunku.
    WYJAŚNIENIE
Twierdzenie 18.
    Jeśli prosta jest prostopadła do dowolnej płaszczyzny, to wszystkie płaszczyzny, które ją zawierają są także prostopadłe do tej płaszczyzny.
    WYJAŚNIENIE
Twierdzenie 19.
    Jeśli dwie płaszczyzny, które się przecinają są prostopadłe do dowolnej płaszczyzny, to ich przecięcie jest również prostopadłe do tej samej płaszczyzny.
    WYJAŚNIENIE
Twierdzenie 20.
    Jeśli kąt trójścienny jest zawarty między trzema kątami płaskimi, to suma dowolnych dwóch z nich jest większa niż pozostały.
    WYJAŚNIENIE
Twierdzenie 21.
    Każdy kąt trójścienny jest zawarty między kątami płaskimi, których suma jest mniejsza od czterech kątów prostych.
    WYJAŚNIENIE
Twierdzenie 22.
    Jeśli dane są trzy kąty płaskie takie, że suma dowolnych dwóch jest większa niż pozostały, i zawarte są one między równymi odcinkami, to można skonstruować trójkąt z trzech odcinków łączących końce ramion tych kątów.
    WYJAŚNIENIE
Twierdzenie 23.
    Skonstruować kąt trójścienny z trzech kątów płaskich, dla których suma dwóch z nich jest większa niż trzeci, tak, aby suma wszystkich trzech kątów była mniejsza od czterech kątów prostych.
    WYJAŚNIENIE
Twierdzenie 24.
    Jeśli bryła jest zawarta między równoległymi płaszczyznami, to przeciwległe ściany są sobie równe i równoległoboczne.
    WYJAŚNIENIE
Twierdzenie 25.
    Jeżeli równoległościan przetniemy płaszczyzną równoległą do przeciwległych ścian, to stosunek otrzymanych brył będzie taki jak stosunek ich podstaw.
    WYJAŚNIENIE
Twierdzenie 26.
    Skonstruować kąt trójścienny równy danemu kątowi trójściennemu przy danym punkcie na prostej.
    WYJAŚNIENIE
Twierdzenie 27.
    Na danej prostej wykreślić równoległościan jednakowo położony i podobny do danego.
    WYJAŚNIENIE
Twierdzenie 28.
    Jeśli równoległościan przecięty jest płaszczyzną poprowadzoną przez przekątne ścian przeciwległych, to jest przez nią podzielony na dwie równe części.
    WYJAŚNIENIE
Twierdzenie 29.
    Równoległościany o tej samej podstawie i wysokości, i których końce krawędzi bocznych leżą na tych samych prostych, są równe.
    WYJAŚNIENIE
Twierdzenie 30.
    Równoległościany, które mają tę samą podstawę i wysokość i, których końce krawędzi bocznych nie leżą na tych samych prostych, są równe.
    WYJAŚNIENIE
Twierdzenie 31.
    Równoległościany, które mają jednakowe podstawy i taką samą wysokość, są sobie równe.
    WYJAŚNIENIE
Twierdzenie 32.
    Stosunek równoległościanów, które mają taką samą wysokość jest taki sam, jak stosunek ich podstaw.
    WYJAŚNIENIE
Twierdzenie 33.
    Stosunek równoległościanów podobnych jest sześcianem stosunku ich odpowiednich krawędzi.
    WYJAŚNIENIE
Twierdzenie 34.
    W równych równoległościanach podstawy są wzajemnie proporcjonalne do wysokości; a równoległościany, w których podstawy są wzajemnie proporcjonalne do wysokości, są równe. 
    WYJAŚNIENIE
Twierdzenie 35.
    Jeżeli dane są dwa równe kąty płaskie, i przez ich wierzchołki poprowadzimy proste nachylone do ich płaszczyzny pod pewnym kątem, tworzące równe kąty z ramionami danych kątów, odpowiednio, jeśli na tych prostych obierzemy dowolnie punkty i zrzutujemy je prostopadle do płaszczyzn, w których leżą dane kąty płaskie, i jeżeli przez otrzymane rzuty punktów na płaszczyźnie poprowadzimy proste do wierzchołków danych kątów, to proste te utworzą z prostymi nachylonymi do płaszczyzny kątów równe kąty.
    WYJAŚNIENIE
Twierdzenie 36.
    Jeśli trzy odcinki są proporcjonalne, to równoległościan utworzony z tych trzech odcinków równy jest równoległościanowi utworzonemu ze średniego z nich, ma on wszystkie ściany równoboczne, ale kąty ma równe kątom wyżej wspomnianego równoległościanu.
    WYJAŚNIENIE
Twierdzenie 37.
    Jeśli cztery odcinki są proporcjonalne, to równoległościany utworzone na nich podobne i podobnie wykreślone są także proporcjonalne; oraz, jeśli równoległościany są podobne i podobnie wykreślone na czterech odcinkach, to odcinki te są również proporcjonalne. 
    WYJAŚNIENIE
Twierdzenie 38.
    Jeśli boki przeciwległych ścian sześcianu przetniemy na dwie równe części i przez punkty przecięcia poprowadzimy płaszczyzny, to przecięcie tych płaszczyzn i przekątna sześcianu przetną się wzajemnie na dwie równe części.
    WYJAŚNIENIE
Twierdzenie 39.
    Jeśli dane są dwa graniastosłupy o równych wysokościach i jeden ma w podstawie trójkąt, a drugi równoległobok oraz jeśli równoległobok jest dwa razy większy od trójkąta (jest jego dwukrotnością), to graniastosłupy te są równe.


\section{Księga dwunasta, o Iksach Dupiksach}

Twierdzenie 1
    Stosunek pól wielokątów podobnych wpisanych w koło jest równy stosunkowi kwadratów średnic tych kół.
    WYJAŚNIENIE
Twierdzenie 2
    Stosunek pól kół jest równy stosunkowi kwadratów długości ich średnic.
    WYJAŚNIENIE
Twierdzenie 3
    Każdy ostrosłup trójkątny można podzielić na dwa równe ostrosłupy trójkątne podobne do danego i na dwa równe graniastosłupy większe niż połowa danego ostrosłupa.
    WYJAŚNIENIE
Twierdzenie 4
    Dane są dwa ostrosłupy trójkątne o takiej samej wysokości i każdy z nich został podzielony na dwa równe ostrosłupy, podobne do całego ostrosłupa oraz na dwa równe graniastosłupy. Podział taki dokonywany jest wielokrotnie. Wtedy stosunek pola podstawy jednego ostrosłupa do pola podstawy drugiego ostrosłupa jest równy stosunkowi graniastosłupów w jednym ostrosłupie do odpowiednich graniastosłupów w drugim ostrosłupie.
    WYJAŚNIENIE
Twierdzenie 5
    Stosunek objętości ostrosłupów trójkątnych o takich samych wysokościach jest równy stosunkowi pól ich podstaw.
    WYJAŚNIENIE
Twierdzenie 6
    Stosunek objętości ostrosłupów o jednakowych wysokościach a podstawach będących dowolnym wielokątem jest równy stosunkowi pól ich podstaw.
    WYJAŚNIENIE
Twierdzenie 7
    Każdy graniastosłup trójkątny można podzielić na trzy ostrosłupy trójkątne równej objętości.
    WYJAŚNIENIE
Twierdzenie 8
    Stosunek objętości ostrosłupów trójkątnych podobnych jest równy sześcianowi stosunku ich odpowiednich krawędzi.
    WYJAŚNIENIE
Twierdzenie 9
    W ostrosłupach trójkątnych o jednakowych objętościach pola podstaw są odwrotnie proporcjonalne do wysokości. Jeżeli ostrosłupy trójkątne mają podstawy odwrotnie proporcjonalne do wysokości, to ich objętości są równe.
    WYJAŚNIENIE
Twierdzenie 10
    Objętość stożka jest jedną trzecią objętości walca o takiej samej podstawie i wysokości.
    WYJAŚNIENIE
Twierdzenie 11
    Stosunek objętości stożka i walca o takiej samej wysokości jest równy stosunków pól ich podstaw.
    WYJAŚNIENIE
Twierdzenie 12
    Stosunek stożków (walców) podobnych jest równy sześcianowi stosunku średnic ich podstaw.
    WYJAŚNIENIE
Twierdzenie 13
    Jeżeli poprowadzimy płaszczyznę równoległą do podstaw walca, to otrzymane dwa walce są do siebie w takim stosunku jak ich wysokości.
    WYJAŚNIENIE
Twierdzenie 14
    Stosunek objętości walców (stożków) o jednakowych podstawach jest równy stosunkowi ich wysokości.
    WYJAŚNIENIE
Twierdzenie 15
    Jeżeli walce lub stożki mają równe objętości to ich wysokości są odwrotnie proporcjonalne do pól podstaw; a jeśli pola podstaw walców (stożków) są odwrotnie proporcjonalne do wysokości, to te walce (stożki) mają jednakowe objętości.
    WYJAŚNIENIE
Twierdzenie 16
    Mając dane dwa koła współśrodkowe należy wpisać w większe koło wielokąt foremny o parzystej liczbie boków tak, aby nie miał on punktów wspólnych z mniejszym kołem.
    WYJAŚNIENIE
Twierdzenie 17
    Dane są dwie kule współśrodkowe. Należy wpisać wielościan w większą kulę tak, aby nie miał on punktów wspólnych z mniejszą kulą.
    WYJAŚNIENIE
Twierdzenie 18
    Stosunek objętości dwóch kul jest równy stosunkowi sześcianów ich średnic.
    WYJAŚNIENIE



    Twierdzenie 1
    Jeżeli odcinek przetniemy w sposób złoty, wtedy kwadrat zbudowany na większej jego części dodany do połowy całego odcinka jest pięciokrotnością kwadratu zbudowanego na połowie tego odcinka.
    WYJASNIENIE
Twierdzenie 2
    Jeżeli pole kwadratu zbudowanego na odcinku jest pięciokrotnością pola kwadratu zbudowanego na jego części, wtedy połowa wspomnianego odcinka jest przecięty w połowie średniego współczynnika, większy segment zaś jest pozostałą częścią pierwotnej prostej.
    WYJASNIENIE
Twierdzenie 3
    Jeżeli linia prosta jest przecięta w połowie średniego współczynnika, to kwadrat zbudowany na sumie krótszego odcinka i połowie większego jest pięciokrotnością kwadratu zbudowanego na połowie dłuższego odcinka.
    WYJASNIENIE
Twierdzenie 4
    Jeżeli linia prosta jest przecięta w połowie średniego współczynnika to suma kwadratów zbudowanych na całości na krótszym odcinku jest trzykrotnością kwadratu na dłuższym odcinku.
    WYJASNIENIE
Twierdzenie 5
    Jeżeli linia jest przecięta w połowie średniego współczynnika, i linia prosta równa dłuższemu odcinkowi jest do niej dodana, wtedy całą linia prosta zostaje przecięta w połowie średniego współczynnika i pierwotna linia prosta jest dłuższym odcinkiem.
    WYJASNIENIE
Twierdzenie 6
    Jeżeli racjonalna linia prosta jest przecięta w połowie średniego współczynnika to każdy z odcinków jest nieracjonalną linią prostą nazywaną apotemą.
    WYJASNIENIE
Twierdzenie 7
    Jeśli trzy kąty pięciokąta foremnego, wzięte w kolejności lub bez niej, wtedy pięciokąt jest równokątny.
    WYJASNIENIE
Twierdzenie 8
    Jeżeli w pięciokącie foremnym weźmiemy dwa odcinki łączące co drugi kąt w uporządkowany sposób, to te odcinki przetną się w sposób złoty, i większe ich części są równe bokowi pięciokąta.
    WYJASNIENIE
Twierdzenie 9
    Jeżeli boki ośmiokąta i sześciokąta wpisanego w ten sam okrąg dodamy do siebie, wtedy cała linia prosta jest przecięta w sposób złoty, i jej więsza część jest bokiem sześciokąta.
    WYJASNIENIE
Twierdzenie 10
    Jeżeli pięciokąt foremny jest wpisany w okrąg, wtedy kwadrat zbudowany na boku pięciokąta równa się sumie kwadratów na bokach sześciokąta i dziesięciokąta wpisanych w ten sam okrąg.
    WYJASNIENIE
Twierdzenie 11
    Jeżeli foremny pięciokąt jest wpisany w okrąg, który ma racjonlaną średnicę, to bok pięciokąta jest irracjonalną linią prostą zwaną minor.
    WYJASNIENIE
Twierdzenie 12
    Jeżeli trójkąt równoboczny jest wpisany w koło, to kwadrat na boku trójkąta jest trzykrotnością kwadratu na promieniu okręgu.
    WYJASNIENIE
Twierdzenie 13
    O konstrukcji piramidy wpisanej w kulę i dowodzie faktu, że kwadrat zbudowany na średnicy jest płótora raza większy od kwadratu zbudowanego na boku piramidy.
    WYJASNIENIE
Twierdzenie 14
    O skonstruowaniu ośmiokąta i wpisaniu do w kulę, tak jak w poprzednim przypadku, i dla udowodnienia, że kwadrat zbudowany na kuli jest dwukrotnością kwadratu zbudowanego na boku ośmiokąta.
    WYJASNIENIE
Twierdzenie 15
    Dla skonstruowania sześcianu wpisanego w kulę, podobnie jak w przypadku piramidy, i dla udowodnienia, że kwadrat zbudowany na średnicy kuli jest trzykrotnością zbudowanego na boku sześcianu.
    WYJASNIENIE
Twierdzenie 16
    Dla skonstruowania dwudziestościanu i wpisania go w kulę, jak w przypadku poprzednich brył, i dla udowodnienia, że kwadrat zbudowany na boku dwudziestościanu jest linią irracjonalną, zwaną minorem.
    WYJASNIENIE
Twierdzenie 17
    Dla skonstruowania dwunastościanu i wpisania go w kulę, jak w przypadku poprzednich figur, i dla udowodnienia, że kwadrat zbudowany na boku dwunastościanu i jest irracjonalną linią zwaną apotomą.
    WYJASNIENIE
Twierdzenie 18
    Dla ustalenia boków pięciu figur i ich porównania.











    (CHYBA 4)
    
Księga VI - Definicje:

Definicja 1.
    Mówi się że figura prostokreślna wpisuje się w figurę prostokreślną, wtedy kiedy każdy kąt figury wpisanej dotyka się każdego boku figury, w który się wpisuje. 
Definicja 2.
    Podobnie się mówi, że figura opisuje się na figurze, kiedy każdy bok figury opisanej dotyka każdego kąta figury na której się opisuje. 
Definicja 3.
    Figura prostokreślna wpisuje się w koło, kiedy każdy kąt figury wpisanej dotyka okręgu koła. 
Definicja 4.
    Figura prostokreślna opisuje się na kole kiedy każdy bok figury opisanej dotyka okręgu koła. 
Definicja 5.
    Podobnież koło wpisuje się w figurę prostokreślną, kiedy każdy bok figury w którą koło się wpisuje, dotyka okręgu koła. 
Definicja 6.
    Koło opisuje się na figurze prostokreślne wtedy gdy okrąg dotyka do każdego kąta figury na której opisujemy koło. 
Definicja 7.
    Mówi się, że linie proste kreśli się w kole, gdy jej końce są na okręgu danego koła. 


Księga VI - Twierdzenia:

Twierdzenie 1.
    W danym kole wykreślić linię prostą równą danej linii prostej, która nie jest większa od średnicy koła.
Twierdzenie 2.
    W dane koło wpisać trójkąt równoramienny względem danego trójkąt.
Twierdzenie 3.
    Na danym kole opisać trójkąt równokątny względem danego trójkąta.
Twierdzenie 4.
    W dany trójkąt wpisać koło.
Twierdzenie 5.
    Na danym trójkącie opisać koło.
Twierdzenie 6.
    W dane koło wpisać kwadrat.
Twierdzenie 7.
    Na danym kole opisać kwadrat.
Twierdzenie 8.
    W dany kwadrat wpisać koło.
Twierdzenie 9.
    Na danym kwadracie opisać koło.
Twierdzenie 10.
    Wykreślić trójkąt równoramienny, którego każdy kąt przy podstawie byłby podwojeniem kąta trzeciego.
Twierdzenie 11.
    W dane koło wpisać pięciokąt równoboczny i równokątny.
Twierdzenie 12.
    Na danym kole opisać pięciokąt równoboczny i równokątny.
Twierdzenie 13.
    W dany pięciokąt równoboczny i równokątny wpisać okrąg.
Twierdzenie 14.
    Na danym pięciokącie równoboczny i równokątny opisać koło.
Twierdzenie 15.
    W dane koło wpisać sześciokąt równoboczny i równokątny.
Twierdzenie 16.
    W dane koło wpisać piętnastokąt równoboczny. 

    Copyright by Bronisław Pabich 2002 - 2007